%% 
%% Copyright 2007-2024 Elsevier Ltd
%% 
%% This file is part of the 'Elsarticle Bundle'.
%% ---------------------------------------------
%% 
%% It may be distributed under the conditions of the LaTeX Project Public
%% License, either version 1.3 of this license or (at your option) any
%% later version.  The latest version of this license is in
%%    http://www.latex-project.org/lppl.txt
%% and version 1.3 or later is part of all distributions of LaTeX
%% version 1999/12/01 or later.
%% 
%% The list of all files belonging to the 'Elsarticle Bundle' is
%% given in the file `manifest.txt'.
%% 
%% Template article for Elsevier's document class `elsarticle'
%% with numbered style bibliographic references
%% SP 2008/03/01
%% $Id: elsarticle-template-num.tex 249 2024-04-06 10:51:24Z rii $
%%
\documentclass[preprint,12pt]{elsarticle}

%% Use the option review to obtain double line spacing
%% \documentclass[authoryear,preprint,review,12pt]{elsarticle}

%% Use the options 1p,twocolumn; 3p; 3p,twocolumn; 5p; or 5p,twocolumn
%% for a journal layout:
%% \documentclass[final,1p,times]{elsarticle}
%% \documentclass[final,1p,times,twocolumn]{elsarticle}
%% \documentclass[final,3p,times]{elsarticle}
%% \documentclass[final,3p,times,twocolumn]{elsarticle}
%% \documentclass[final,5p,times]{elsarticle}
%% \documentclass[final,5p,times,twocolumn]{elsarticle}

%% For including figures, graphicx.sty has been loaded in
%% elsarticle.cls. If you prefer to use the old commands
%% please give \usepackage{epsfig}

%% The amssymb package provides various useful mathematical symbols
\usepackage{amssymb}
%% The amsmath package provides various useful equation environments.
\usepackage{amsmath}
\usepackage{lyu}
\usepackage{algorithm}
\usepackage{algpseudocode}%[noend]
%% The amsthm package provides extended theorem environments
\usepackage{amsthm}\mathchardef\mhyphen="2D
% \usepackage{amsmath,amsthm,amssymb}
\newtheorem{theorem}{Theorem}
\newtheorem{lemma}[theorem]{Lemma}
\theoremstyle{definition}
\newtheorem{definition}[theorem]{Definition}
\theoremstyle{remark}
\newtheorem{remark}[theorem]{Remark}
\newtheorem{proposition}{Proposition}
\newtheorem{assumption}{Assumption}
\newtheorem{corollary}{Corollary}
%% The lineno packages adds line numbers. Start line numbering with
%% \begin{linenumbers}, end it with \end{linenumbers}. Or switch it on
%% for the whole article with \linenumbers.
%% \usepackage{lineno}
\journal{Journal of Computational Physics}
% \journal{Nuclear Physics B}

\begin{document}

\begin{frontmatter}

%% Title, authors and addresses

%% use the tnoteref command within \title for footnotes;
%% use the tnotetext command for theassociated footnote;
%% use the fnref command within \author or \affiliation for footnotes;
%% use the fntext command for theassociated footnote;
%% use the corref command within \author for corresponding author footnotes;
%% use the cortext command for theassociated footnote;
%% use the ead command for the email address,
%% and the form \ead[url] for the home page:
%% \title{Title\tnoteref{label1}}
%% \tnotetext[label1]{}
%% \author{Name\corref{cor1}\fnref{label2}}
%% \ead{email address}
%% \ead[url]{home page}
%% \fntext[label2]{}
%% \cortext[cor1]{}
%% \affiliation{organization={},
%%             addressline={},
%%             city={},
%%             postcode={},
%%             state={},
%%             country={}}
%% \fntext[label3]{}

\title{Stochastic Branching method for solving the Keller--Segel chemotaxis Model}

%% use optional labels to link authors explicitly to addresses:
%% \author[label1,label2]{}
%% \affiliation[label1]{organization={},
%%             addressline={},
%%             city={},
%%             postcode={},
%%             state={},
%%             country={}}
%%
%% \affiliation[label2]{organization={},
%%             addressline={},
%%             city={},
%%             postcode={},
%%             state={},
%%             country={}}

\author[1,2]{Liyao Lyu}
\ead{lyuliyao@math.ucla.edu}
\author[2,3]{Huan Lei}
\ead{leihuan@msu.edu}

\address[1]{{Department of Mathematics, University of California, Los Angeles}, 520 Portola Plaza, Los Angeles, CA, USA}

\address[2]{{Department of Computational Mathematics, Science \& Engineering, Michigan State University}, 428 S Shaw Ln, East Lansing, MI, USA}

\address[3]{{Department of Statistics \& Probability, Michigan State University}, 619 Red Cedar Road, East Lansing, MI, USA}


%% Abstract
\begin{abstract}
We propose an efficient stochastic branching method for solving parabolic-parabolic Keller--Segel (KS) chemotaxis systems. In this approach, both the cell density and the chemical concentration are modeled as empirical measures of particles. By approximating the original problem linearly through a Kolmogorov forward equation for the particles at each time step, our method provides an effective solution to a system of advection--diffusion--reaction equations. Compared to traditional methods, our approach offers better handling of singularities. We present numerical examples in two-dimensional space to demonstrate the accuracy and effectiveness of the proposed method.
\end{abstract}

%%Graphical abstract
%\begin{graphicalabstract}
%\includegraphics{grabs}
%\end{graphicalabstract}

%%Research highlights
% \begin{highlights}
% \item Research highlight 1
% \item Research highlight 2
% \end{highlights}

%% Keywords
\begin{keyword}
%% keywords here, in the form: keyword \sep keyword

%% PACS codes here, in the form: \PACS code \sep code

%% MSC codes here, in the form: \MSC code \sep code
%% or \MSC[2008] code \sep code (2000 is the default)

\end{keyword}

\end{frontmatter}

%% Add \usepackage{lineno} before \begin{document} and uncomment 
%% following line to enable line numbers
%% \linenumbers

%% main text
%%

%% Use \section commands to start a section
\section{Introduction}
The advection--diffusion--reaction type partial differential equations (PDEs) have attracted much interest in physics, biology, and engineering for modeling various phenomena relevant to chemical reactions, transport, and collective behaviors in biological systems \cite{hundsdorfer2013numerical}. Traditional mesh-based methods, such as finite difference~\cite{hoff1978stability}, finite element~\cite{codina1998comparison}, and discontinuous Galerkin schemes~\cite{ayuso2009discontinuous,xia2009application}, have achieved great success in solving many of these problems. However, from a numerical perspective, an essential challenge in solving such PDEs lies in handling nonlinearities and singular structures in some special cases, e.g., the Keller--Segel (KS) model \cite{keller1970initiation}.  
Mesh-based discretizations approximate solutions in function spaces that are not well-suited for capturing blow-up or highly localized behavior, often resulting in reduced accuracy or numerical instabilities.


Particle-based approaches offer a natural alternative by representing the solutions as empirical measures, which are particularly effective for problems exhibiting concentration or blow-up phenomena. 
Existing particle methods, such as the particle-in-cell method~\cite{birdsall2018plasma, hockney2021computer} for solving the Vlasov--Poisson equation, are often based on well-established  underlying physical laws. In this setting, particle dynamics can be derived from the explicit micro-scale governing equations. In contrast, the explicit physical laws are often unknown for general reaction–diffusion–advection type PDEs, where effective interactions may be highly nonlinear, system-specific, or even not well characterized. For example, in many biological and chemical systems such as chemotaxis, the micro-scale interactions remain poorly understood, which limits the applicability of traditional particle-based approaches.
To extend particle-based methods beyond physically motivated settings, blob methods~\cite{craig2016blob,carrillo2019blob,huang2017error} and semi-Lagrangian particle schemes~\cite{qiu2010conservative,qiu2011conservative,qiu2011positivity} have been developed for conservative PDEs.
% To provide particle-based methods for a more general form of PDEs, recent developments, such as blob methods \cite{craig2016blob,carrillo2019blob,huang2017error} and semi-Lagrangian particle schemes \cite{qiu2010conservative,qiu2011conservative,qiu2011positivity}, have provided accurate and stable discretizations for a variety of conservative PDEs. 
These approaches are based on the mean-field perspective, in which the empirical measures of interacting particle systems converge to the solution of the original PDE. This mathematical framework establishes a rigorous link between the particle dynamics and the continuum equations and does not rely on explicit microscale physical laws. Despite these advances, there is still no unified particle-based framework for systematically discretizing nonlinear and nonconservative PDEs, especially in the presence of singularities or blow-up phenomena.

Machine learning-based numerical PDE algorithms, such as Deep Galerkin Methods~\cite{sirignano2018dgm}, Deep Ritz method~\cite{yu2018deep}, and Mixed Residual method~\cite{lyu2022mim}, have attracted significant attention in recent years. By leveraging the strong representability of neural networks, these methods are particularly promising for solving high-dimensional problems. In the context of particle-based approaches, coupling neural networks with the Feynman-Kac formula~\cite{han2024deep, hutzenthaler2021speed, nguwi2024deep,beck2021deep} or backward stochastic differential equations (SDEs) ~\cite{han2017deep, han2018solving, hure2020deep} has also emerged as a powerful strategy for solving a broad class of PDEs. Nevertheless, the approximation capability of neural networks is fundamentally grounded in the universal approximation theorem that targets continuous solutions on compact sets.
As a result, neural-network–based solvers may face intrinsic difficulties when dealing with PDEs that develop singularities or blow-up phenomena, where discontinuities or concentration effects play a central role.

In this work, we propose a \textit{stochastic branching method} to address these challenges. 
At each time step, the original PDE is locally linearized and approximated by the Kolmogorov forward equation, which establishes a direct connection between the PDE and the stochastic particle dynamics.  Unlike approaches that rely on explicit physical laws or mean-field limit arguments, this framework establishes a systematic particle-based discretization applicable to general nonlinear and nonconservative PDEs. 
In particular, a birth–death mechanism is introduced to account for the non-conservative nature of the system which adjusts particle numbers consistently with the underlying reaction dynamics.  
The only step that invokes a continuity assumption is the interpolation of the empirical measure, which can be replaced by blob- or kernel-based representations. Consequently, the developed method does not fundamentally rely on continuity and is well-suitable for problems with singularities or blow-up phenomena. From a stochastic particle perspective, the algorithm is naturally consistent with the physical intuition behind Keller–Segel dynamics, in line with Stevens’ mean-field derivation of the Keller–Segel equations from moderately interacting particle systems \cite{Stevens_SIAM_2000} and can be extended naturally to a broader class of advection–diffusion–reaction systems.
% formulation is closely related to Stevens’ mean-field derivation of Keller--Segel equations from moderately interacting particle systems \cite{Stevens_SIAM_2000}. Unlike  \cite{Stevens_SIAM_2000}, we develop a finite-$N$ particle discretization based on local Kolmogorov linearization and a particle-level birth-death update, which directly handles the nonlinear and nonconservative terms and extends to a class of advection–diffusion–reaction systems.  


To illustrate the essential idea, we consider the two-dimensional Keller--Segel chemotaxis model~\cite{keller1970initiation, patlak1953random} as the theme problem, which takes the form 
\begin{equation}\label{equ:KS}
\begin{aligned}
\partial_t u - \nabla \cdot ( \nabla u - \chi(u) \nabla v) &= f_u(u,v), & \bx\in \mathbb{R}^2 , t\in(0,T), \\
\partial_t v - \Delta v &= f_v(u,v), & \bx\in \mathbb{R}^2 , t\in(0,T),
\end{aligned}
\end{equation}
where $u$ and $v$ represent the cell density and chemical concentration, respectively. $\chi$ represents the chemotactic sensitivity, and $f_u$ and $f_v$ represent the production or decay rate.  For simplicity, we focus on periodic boundary conditions or unbounded domains, though the framework can be extended to more general boundary settings in future research. 
%Chemotaxis is a process in which cells navigate toward areas with favorable chemical concentrations and away from unfavorable ones. 
The Keller--Segel system~\eqref{equ:KS} has been extensively studied over the past several decades, with numerous works addressing both analytical properties and numerical methods across model variants and dimensions. In two dimensions, it is well understood that both parabolic–elliptic and fully parabolic cases can exhibit finite-time blow-up to a Dirac-delta profile when the initial mass is sufficiently large~\cite{herrero1997blow, perthame2007transport}. Furthermore, in the three-dimensional parabolic–elliptic case, an alternative self-similar blow-up profile of the form $ C(T-t+\vert \mb x\vert ^2)^{-1}$ has also been identified in the literature~\cite{giga2011asymptotic,souplet2019blow}. 

On the numerical side, many notable methods have been developed for the KS systems. The finite-volume method~\cite{filbet2006finite,chertock2008second,chertock2018high} achieves efficient and positivity-preserving simulations for chemotaxis and haptotaxis models.  
%
Shen et al.~\cite{shen2020unconditionally} introduced an energy-dissipative and bound-preserving formulation, subsequently extended by Chen et al.~\cite{chen2022error} to a fully discrete finite element scheme for the 2D parabolic–elliptic KS, which reproduces finite-time blow-up under assumptions analogous to the continuous setting. 
Liu and Wang~\cite{liu2018positivity} reformulated the equations via the Le Ch\^atelier principle to construct positivity-preserving schemes. Discontinuous Galerkin methods have also been applied to KS systems~\cite{zhong2024direct, li2017local,epshteyn2009new}, in addition to various mesh-based finite element, finite volume, and domain-decomposition techniques for 3D fully parabolic models~\cite{epshteyn2019efficient, strehl2013positivity}. For particle-based approaches, the most relevant work is the recently proposed stochastic interacting particle–field algorithm~\cite{wang2025novel,hu2024stochastic,chertock2025hybrid}. However, this approach does not deal with the non-conservative equations at the particle level, limiting its applicability to the broader class of nonlinear nonconservative PDEs. 

We address this challenge by explicitly coupling the SDE-based advection--diffusion with a particle-level birth--death mechanism for reactions, which overcomes the limitations of traditional methods and yields a particle measure discretization that remains robust near concentration and blow-up. We validate the method on two-dimensional KS benchmarks, including near blow-up and nonlinear growth, and compare with the standard grid-based method, observing the expected $1/\sqrt{N}$ scaling with the particle number.
%
The remainder of the paper is organized as follows: Section 2-3 present the stochastic branching method with a heuristic derivation of the limiting dynamics; Section 4 extends the approach for the KS systems; Section 5 reports numerical results and Section 6 discusses extensions and outlook.

% This work makes three contributions. (\Rmnum{1}) We introduce a stochastic branching framework that couples SDE-based advection–diffusion with a birth–death mechanism for reactions, yielding a particle–measure discretization that is not tied to smooth function spaces and remains effective near concentration and blow-up. (\Rmnum{1}) We derive the scheme via local Kolmogorov linearization and an operator splitting, giving a transparent PDE–SDE correspondence and an implementable algorithm with Fourier/kernel reconstructions of empirical measures. (\Rmnum{1}) We validate the method on two-dimensional KS benchmarks—including near-blow-up and nonlinear growth—against standard grid-based baselines, observing the expected $1/\sqrt{N}$ scaling with particle number. We also discuss extensions to other reaction–diffusion systems (Gray–Scott, Allen–Cahn, Schnakenberg) and paths toward higher-dimensional settings, with details deferred to the Discussion


% In summary, we propose a novel stochastic branching method for solving the Keller–Segel chemotaxis model, which overcomes the limitations of traditional numerical methods by efficiently handling singularities and nonlinearities at particle level. A heuristic derivation of the algorithm is presented in the current version and will be discussed in detail in future work. Although our focus here is on the two-dimensional Keller–Segel model, the proposed framework is general and can be applied to a broad class of reaction–diffusion systems, such as the Gray–Scott model, the Allen–Cahn equation, and Schnakenberg pattern formation. Furthermore, by leveraging recent advances in density estimation techniques~\cite{novikov2021tensor, tang2025wavelet, sheng2025approximation, tang2024solving}, our approach can be naturally extended to higher-dimensional systems. The effectiveness of the method is demonstrated through several numerical examples, highlighting both its robustness and potential for generalization.


\section{Stochastic branching method}

Let us start with a scalar non-conservative advection-diffusion–reaction equation
\begin{equation}\label{equ:reaction-diffusion}
  \partial_t u \;=\; \nabla \!\cdot\! \big( a(u)\,u \,+\, \nabla\big(d(u)\,u\big) \big) \;+\; r(u)\,u
\end{equation}
where $a:\mathbb{R}_{+}\!\to\!\mathbb{R}^2$ is the advective velocity, $d:\mathbb{R}_{+}\!\to\!\mathbb{R}_{+}$ is the diffusion coefficient, and $r:\mathbb{R}_{+}\!\to\!\mathbb{R}$ is the reaction rate. The extension to the full KS system~\eqref{equ:KS} is presented in Section~\ref{sec:extension}.



Let $S(N,t)=\{\bX_N^i\}_{i\in\mathcal I(N,t)}$ denote a set of particles at time $t$ whose initial population is $N$, where the particles' indices are numbered consecutively in the index set $\mathcal I(N,t)$ and $\bX_N^i(t)$ denote the position of the $i\mhyphen$th particle. The associated empirical (finite) measure is
\begin{equation}
    t \to \mu^N(t,\intd \bx) =  \frac{1}{N}\sum_{i\in \mathcal I(N,t)}\delta_{\bX_N^i(t)},
\end{equation}
where $\delta_\bx $ denotes the Dirac measure at $\bx\in \mathbb {R}^2$. 

The core idea of particle-based methods is to approximate the PDE solution by the limit of the empirical measure
\[
\mu(t,\mathrm{d}\bx) = \lim_{N\to\infty} \mu^N(t,\mathrm{d}\bx).
\]
Since Eq. \ref{equ:reaction-diffusion} is nonconservative, the total mass of the system is not preserved. Consequently, both $\mu^N$ and its limit $\mu$ should be viewed as finite measures rather than probability measures.


\subsection{Initial condition and Normalization}
To connect the solution $u(t,\bx)$ to the measure $\mu(t,\mathrm{d}\bx)=\rho(t,\bx) \intd \bx$, we introduce a linear transformation between them. 
Without loss of generality, we assume that the initial measure $\mu(0,\cdot)$ is a probability measure with total mass $1$, and define
\begin{equation*}
    u(t,\bx)\,\mathrm{d}\bx = Z \,\mu(t,\mathrm{d}\bx),
\end{equation*}
where 
\[
Z = \int_{\Omega} u(0,\bx)\,\mathrm{d}\bx
\]
denotes the total initial mass.
Thus $\rho(0,\cdot)$ is a probability density
\[
\rho(0,\mathbf{x}) \;=\; \frac{u(0,\mathbf{x})}{Z},
\]
and we can sample the initial particle positions independently from $\rho(0,\cdot)$ (e.g., via Metropolis–Hastings) to form $S(N,0)=\{\mathbf{X}^i_0\}_{i=1}^N$. 
Substituting $u=Z\rho$ into~\eqref{equ:reaction-diffusion}  yields
\begin{equation}\label{equ:normalized_diffusion_reaction}
\partial_t \rho 
= \nabla \cdot \left( \tilde a(\rho)\rho + \nabla \big(\tilde d(\rho)\rho \big)\right) 
  + \tilde r(\rho)\rho,
\end{equation}
with rescaled coefficients
\[
\tilde a(\rho) = a(Z\rho), 
\qquad \tilde d(\rho) = d(Z\rho), 
\qquad \tilde r(\rho) = r(Z\rho).
\]
Since reactions generally change the total mass, $\int_{\Omega}\rho(t,\mathbf{x})\,\mathrm{d}\mathbf{x}$ may not equal $1$ for $t>0$. Rather, it tracks the mass of $u$ via $\int_{\Omega} u(t,\mathbf{x})\,\mathrm{d}\mathbf{x} = Z \int_{\Omega} \rho(t,\mathbf{x})\,\mathrm{d}\mathbf{x}$.

\subsection{Operator Splitting}
To approximate the solution of the PDE, we discretize the time interval $[0,T]$ with step size $\tau>0$ and grid points
\[
0 = t_0 < t_1 < \cdots < t_M = T,\qquad \tau=t_{n+1}-t_n.
\]
We then employ an operator-splitting strategy that separates the dynamics according to their physical mechanisms. 
Equation~\eqref{equ:normalized_diffusion_reaction} can be written as
\begin{equation}
\partial_t \rho = \mathcal{A}\rho + \mathcal{B}\rho ,
\end{equation}
where the operators $\mathcal{A}$ and $\mathcal{B}$ correspond to
\begin{align}
\mathcal{A}\rho &= \nabla \cdot \left( \tilde a(\rho)\rho + \nabla \!\big(\tilde d(\rho)\rho \big) \right), \\[4pt]
\mathcal{B}\rho &= \tilde r(\rho)\rho .
\end{align}
Here, $\mathcal{A}$ describes the advection--diffusion dynamics, while $\mathcal{B}$ accounts for the reaction mechanism (e.g., proliferation or death).  The evolution of the exact solution over $t$ can be expressed formally as
\begin{equation}
\rho(t,\bx) = \exp\left( t (\mathcal{A} + \mathcal{B}) \right)\rho(0,\bx).
\end{equation}
In practice, the operator $\exp\!\big(\tau(\mathcal{A}+\mathcal{B})\big)$ cannot be evaluated explicitly, 
we adopt the first-order Lie-Trotter splitting \cite{trotter1959product}, which is justified by the Trotter product formula:
\begin{equation}
\exp\!\left( t (\mathcal{A} + \mathcal{B}) \right) 
= \lim_{M\to\infty} \Big( \exp\!\left(\tfrac{t}{M} \mathcal{B}\right) 
                          \exp\!\left(\tfrac{t}{M} \mathcal{A}\right) \Big)^M .
\end{equation}
At the discrete level, we approximate $\rho(t_{n+1},\cdot)$ by
\begin{equation}
\rho_{n+1} = \exp(\tau \mathcal{B}) \exp(\tau \mathcal{A}) \rho_n ,
\end{equation}
for $n=0,\cdots,M-1$.
This scheme provides a simple and efficient first-order approximation; 
higher-order splittings, such as Strang splitting \cite{strang1968construction}, may be employed for improved accuracy.
In practice, this decomposition allows us to design particle-based updates for each subprocess separately:
\begin{equation}
\rho_n^\ast = \exp(\tau \mathcal{A})\rho_n, 
\qquad 
\rho_{n+1} = \exp(\tau \mathcal{B})\rho_n^\ast.
\end{equation}

\subsection{Advection Diffusion Step}
\label{sec:adevection_diffusion}
In this step, we propagate particles under the advection–-diffusion dynamics, where each particle follows a stochastic path determined by both a drift and a diffusion term. At time step $n$, let $S^N_n = \{\bX^i_{N,n}\}_{i\in \mathcal I_n^N}$ denote  a set of particles  whose initial population is $N$, which induce the empirical measure 
\[
\mu^N_{n}= \frac{1}{N}\sum_{i\in \mathcal I_n^N}\delta_{\bX_{N,n}^i}
\] 
approximates the solution at step $n$, i.e.,  $\mu^N_{n} (\intd \bx) \approx \mu_n(\intd \bx) = \rho_n(\bx)\intd \bx$. 
Our goal is to evolve these particles according to the transition probability $p(\bx,s;\by,t)$, which leads to an updated empirical measure
\[\mu_n^{N,*} = \int p(\bx,t_n;\by,t_n+\tau)  \mu_n^N(\intd \bx) =   \frac{1}{N}\sum_{i\in \mathcal I_n^N}\delta_{\bX_{N,n}^{i,*}}
\]  
reflecting the action of the advection--diffusion operator on the density.
By the Kolmogorov forward equation, the transition density $p(\bx,s;\by,t)$ satisfies 
\begin{equation}
    \partial_t p = \tilde {\mathcal A }p := - \nabla \cdot(b  p) + \frac{1}{2}\nabla^2:(\sigma \sigma^\top p) 
\end{equation}
with initial condition $\lim_{t\to s^+} p(\bx,s;\cdot,t) = \delta_\bx$. By matching the differential operator $\mathcal A$ with the target PDE \eqref{equ:reaction-diffusion}, we have 
\begin{equation}
    b(x) = -\tilde a(\rho^N_n(x)), \quad \frac{1}{2}\sigma(x) \sigma(x)^\top = \tilde d(\rho^N_n(x)),
\end{equation}
where $\rho^N_n(x)$ represents the interpolated probability density of $\mu_n^N$ and will be discussed in Sec. \ref{sec:interpolation}.
%
Accordingly, the evolution of each particle $\bX^i_n$ as a stochastic process satisfying the It\^o stochastic differential equation
\begin{equation}\label{equ:overd_langevin}
    \intd \bX_{N,n}^i = b(\bX_{N,n}^i) \intd t  + \sigma(\bX_{N,n}^i) \intd \bW^i_{N}(t),
\end{equation}
where $\bW^i_{N}(t)$ is a $2\mhyphen$dimensional Brownian motion. 
In practice, we approximate Eq. \eqref{equ:overd_langevin} using an Euler-Maruyama scheme with step size $\tau$, 
\begin{equation}
 \bX_{N,n}^{i,*} = \mathbf{X}_{N,n}^{i} + b(\bX_{N,n}^i) \tau  + \sigma(\bX_{N,n}^i) \sqrt{\tau} \boldsymbol{\xi}^i_{N,n},
 \label{eq:euler_maruyama}
\end{equation}
where $\boldsymbol{\xi}^i_{N,n}$ represent identically independent distributed normal random variables; see Algorithm \ref{alg:sde} for the summary of the particle propagation for the advection--diffusion step. 

\begin{algorithm}
\caption{SDE propagation of the advection--diffusion step}\label{alg:sde}
\begin{algorithmic}
\Require Particle set: $S=\{\mathbf{X}^i\}_{i\in\mathcal I}$; time step $\tau>0$
\Require Drift term $b(\cdot)$
\Require Diffusion operator $\sigma\in\mathbb R^{2 \times 2}$
\Ensure Propagated set $S_*=\{\mathbf{X}^{i,*}\}_{i\in\mathcal I}$
\Procedure{SDE\_Propagate}{$S,\,\tau,\,b,\,\sigma$}
  \State $S_{*}\gets [\ ]$
  \For{$i\in\mathcal I$}
    \State Sample $\boldsymbol{\xi}_i \sim \mathcal N(0,I_d)$ independently
    \State $\mathbf{X}^{i,*} \gets \mathbf{X}^i \;+\; b(\mathbf{X}^i;\mathcal F)\,\tau \;+\; \sqrt{\tau}\,\sigma\,\boldsymbol{\xi}_i$
    \State Append $\mathbf{X}^{i,*}$ to $S_{*}$
  \EndFor
  \State \Return $S_{*}$
\EndProcedure
\end{algorithmic}
\end{algorithm}

\subsection{Branching Step}
\label{sec:branching}
This step modifies the mass of the empirical measure and change the particle number. Given the intermediate empirical measure from the advection--diffusion step $\mu_{N,n}^*$, we aim to evolve this measure under the action of the reaction operator $\mathcal  B$ over a time step $\tau$ following
\[
\partial_t \rho = c_n(\mathbf{x})\,\rho,
\qquad c_n(\mathbf{x}) := \tilde r\!\big(\rho_n^\ast(\mathbf{x})\big),
\]
such that the resulting empirical measure
\begin{equation}
    \mu^N_{n+1} = \frac{1}{N} \sum_{i\in \mathcal I (N,(n+1)\tau)} \delta_{\bX^i_{N,n+1}}
\end{equation}
approximates the updated density $\mu^N_{n+1}(\intd \bx ) = \rho^N_{n+1}(\bx) \intd \bx \approx {\rm e}^{c_n(\mathbf{x}) \tau} \rho_{n}^{N,*}(\bx) \intd \bx$. Therefore, the particles located at $\bx$ should birth and death with probability $c(\bx)=\tilde r(\rho^*(\bx))$ depending on the positiveness of $\tilde r$.
To implement this step, we adopt a branching process. For each particle $\bX^{i,\ast}_{n}$, we associate a branching rule determined by $c_i:=c_n(\mathbf X^{i,\ast}_{N,n})$. Specifically, over the interval $[t_n,t_{n+1}]$, each particle undergoes one of the following stochastic events:
\begin{itemize}
  \item \emph{If \(c_i\ge 0\) (birth):} keep the parent and draw \(K_i\sim \mathrm{Poisson}\!\big(e^{c_i\tau}-1\big)\); output \(1+K_i\) identical copies at \(\mathbf X^{i,\ast}_{N,n}\).
  \item \emph{If \(c_i<0\) (death):} keep the parent with probability \(p_i=e^{c_i\tau}\), otherwise remove it.
\end{itemize}
% \begin{itemize}
%     \item Survival (no branching): particle $\bX^{i,*}_{N,n}$ survives with probability approximately $\exp(-| c(\bX^{i,*}_{N,n})|\tau)$.
%     \item Death (removal): if $c<0 $, particle $\bX^{i,*}_{N,n}$ dead  with the probability $1-\exp( c(\bX^{i,*}_{N,n})\tau)$.
%     \item Birth (duplication ): if $ c>0 $, particle $\bX^{i,*}_{N,n}$ duplication with the probability $1-\exp(- c (\bX^{i,*}_{N,n})\tau)$.
% \end{itemize}

Let the resulting set be \(S_{n+1}^N=\{\mathbf X^i_{N,n+1}\}_{i\in\mathcal I_{n+1}^N}\) and
\[
\mu_{n+1}^{N}(\mathrm{d}\mathbf{x})
\;=\;
\frac{1}{N}\sum_{i\in\mathcal I_{n+1}^N}\delta_{\mathbf X^i_{N,n+1}}(\mathrm{d}\mathbf{x}).
\]


\begin{proposition}
    Let $S=\{\bX^i\}_{i\in \mathcal I}$ be a set of particles, whose empirical measure is defined as $\mu (\intd \bx) = \frac{1}{N}\sum_{i\in \mathcal I }\delta_{\bX^i}$. If the particles are updated according to the above rules then the updated particle set $S_+=\{ \bX^i_+\}_{i\in {\mathcal I}_+}$ and corresponding  measure $ \mu_+ (\intd \bx) = \frac{1}{N}\sum_{i\in  {\mathcal  I}_+}\delta_{ \bX^i_+}$ should satisfies 
    \begin{equation}
        \mathbb E [\left<f, \mu_+\right>] =   \left<f,\exp( c\tau)\mu\right>,
    \end{equation}
    for all $f\in C^2_b(\mathbb R)$.
\label{prop:stochastic_branch}
\end{proposition}

\begin{proof} Let $\lambda_i = {\rm e}^{c(\mb X^i) \tau} -1 $ denote the  rate of the Poisson process for the birth branch with $c(\mb x^i) \ge 0$. For any test function $f\in C^2_b(\mathbb R^+, \mathbb R^2)$, we have
\begin{equation}
\begin{aligned}
     & \mathbb E\left[\left<f,  \mu_+\right>\right]
    \\=& \frac{1}{N} \sum_{i\in \mathcal I ,c(\bX^i)>0} \sum_{j=0}^{\infty} 
    \left[(j+1) \frac{e^{-\lambda_i} \lambda_i^j}{j!}\right] f(\bX^i)
     + \frac{1}{N} \sum_{i\in \mathcal I , c(\bX^i)<0}
    \exp( c(\bX^i)\tau)f(\bX^i) \\
    % \\=& \frac{1}{N} \sum_{i\in \mathcal I ,c(\bX^i)>0} 
    % \left[\exp(- c(\bX^i)\tau)f(\bX^i) + (1-\exp( -c(\bX^i)\tau))2f(\bX^i)\right]\\
    % & + \frac{1}{N} \sum_{i\in \mathcal I , c(\bX^i)<0}
    % \exp( c(\bX^i)\tau)f(\bX^i) \\
    =& \frac{1}{N}\sum_{i\in \mathcal I ,c(\bX^i)>0} 
    \left(
    \exp(c(\bX^i)\tau)f(\bX^i)
    \right) 
     + \frac{1}{N}\sum_{i\in \mathcal I , c(\bX^i)<0}\left(
    \exp( c(\bX^i)\tau)f(\bX^i) \right)\\
    = & \frac{1}{N}\sum_{i\in \mathcal I}\exp( c(\bX^i)\tau)f(\bX^i) 
    = \left<f,\exp( c\tau)\mu\right>,
\end{aligned}
\end{equation}
where the second equality follows from the fact that the expectation of a Poisson process equal to the rate.
% where the first equality comes from the expectation number of population after birth is $2$ and death if $0$. 
\end{proof}
By choosing $c(\mb x) = \tilde{r}(\rho^{\ast}(\mb x))$, this stochastic branching process realizes the operator $\exp(\mathcal{B}\tau)$ at the particle level. We summarize the process in Algorithm \ref{alg:birthdeath}.

\begin{algorithm}
\caption{Birth--Death Step}\label{alg:birthdeath}
\begin{algorithmic}
\Require Propagated particles $S=\{\mathbf X^{i}\}_{i\in\mathcal I}$, time step $\tau>0$, rate function $r:\mathbb R^2\to\mathbb R$
\Ensure Post-branching set $S_{+}=\{\mathbf X^{i}\}_{i\in\mathcal I_{+}}$
\Procedure{Birth\_Death}{$S,\,\tau,\,r(\cdot)$}
    \State $S_{n+1}\gets [\ ]$
    \For{$i\in\mathcal I$}
        \State $r_i \gets r(\mathbf X^{i})$
        \If{$r_i>0$} \Comment{birth}
            \State Append $\mathbf X^{i}$ to $S_{+}$ \Comment{keep parent}
            \State $p_{\mathrm{birth}} \gets \exp(r_i\,\tau) - 1$
            %\State With probability $p_{\mathrm{birth}}$, append a duplicate of $\mathbf X^{i}$ to $S_{+}$
            \State Draw $K_i \sim \mathrm{Poisson}\!\big(p_{\mathrm{birth}}\big)$, append $K_i$ duplicates of $\mathbf X^{i}$ to $S_{+}$
        \ElsIf{$r_i<0$} \Comment{death}
            \State $p_{\mathrm{survive}}\gets \exp(r_i\,\tau)$ 
            \State With probability $p_{\mathrm{survive}}$, append $\mathbf X^{i}$ to $S_{+}$
            \State \quad (otherwise the particle is removed)
        \Else \Comment{$r_i=0$}
            \State Append $\mathbf X^{i}$ to $S_{+}$
        \EndIf
    \EndFor
    \State \Return $S_{+}$
\EndProcedure
\end{algorithmic}
\end{algorithm}


\subsection{Nonlinear Interpolation}
\label{sec:interpolation}
To evaluate the functions $\tilde a(\rho)$ ,  $\tilde d(\rho)$ , and $\tilde r(\rho)$ , that govern the drift, diffusion, and reaction dynamics in our particle-based scheme, we need to reconstruct the continuous density 
$\rho$ from its empirical measure. Specifically, given an empirical measure
\[
\mu^N_n= \frac{1}{N}\sum_{i\in \mathcal I^N_n}\delta_{\bX_{N,n}^i},
\] 
we seek a smooth approximation of $\rho^N_n$ that allows pointwise evaluation and differentiation as required by the operator $\mathcal{A}$ and $\mathcal{B}$.

In this work, we choose to interpolate the empirical distribution in the Fourier space. Specifically, we project the density onto a finite set of Fourier basis functions, allowing efficient representation and differentiation. The Fourier coefficients are computed from the particle locations $\bX_{N,n}^i$ in periodic boundary condition $\Omega=[-\pi,\pi]^2$ and the reconstructed density 
$\rho_n(t,x)$ is obtained from the Fourier coefficients. To be more specific, we define a Fourier interpolation operator $\mathcal F_K$
that maps the empirical particle distribution into a truncated Fourier representation 
\[
\rho^N_n(\bx) = \mathcal{F}_K(\mu^N_n)(\bx) 
          = \sum_{\vert k_1\vert ,\vert k_2\vert \leq K} \hat{\rho}_n(k_1,k_2)\,\phi_{k_1}(\bx_1)\phi_{k_2}(\bx_2),
\]
where the basis functions are chosen as
\[
\phi_k(x) =
\begin{cases}
\frac{1}{\sqrt{2\pi}}, & k=0, \\
\frac{1}{\sqrt{2\pi}} \cos(k  x), & k \in \mathbb{Z}_+, \\
\frac{1}{\sqrt{2\pi}}\sin(k  x), & k \in \mathbb{Z}_- .
\end{cases}
\]

The Fourier coefficients are computed directly from the empirical measure:
\begin{equation}
\hat{\rho}_n(k_1,k_2) = \int \phi_{k_1}(y_1)\phi_{k_2}(y_2)\,\mu^N_n(\intd\by) 
= \frac{1}{N}\sum_{i\in \mathcal I^N_n} \phi_{k_1}(\bX_{N,    n,1}^{i})\phi_{k_2}(\bX_{N,n,2}^{i}).
\label{eq:Fourier_mode}
\end{equation}
The process is shown in Algorithm \ref{alg:interpolate}. 
This choice is motivated by the simplicity and spectral accuracy of Fourier methods for smooth periodic functions. However, we note that other functional spaces, such as polynomial bases, wavelets, or kernel-based representations, could offer better adaptability in non-periodic or highly localized settings. We leave the exploration of such alternatives to future work.



\begin{algorithm}
\caption{Interpolation Step (fixed-$N$ normalization)}\label{alg:interpolate}
\begin{algorithmic}
\Require Particle set: $S=\{ \mathbf{X}^i \}_{i\in \mathcal{I}}$; fixed initial population $N$; truncation index set $\mathcal{M}_K$; domain descriptor $\Omega$; basis $\{\psi_m\}_{m\in\mathcal{M}_K}$ with (optionally) known gradients $\{\nabla\psi_m\}_{m\in\mathcal{M}_K}$
\Ensure Truncated field $\tilde{\rho}_n(\mathbf{x})$ (and, if needed, $\nabla\tilde{\rho}_n(\mathbf{x})$)
\Procedure{Interpolate}{$S_n,\,N,\,\mathcal{M}_K,\,\Omega,\,\{\psi_m\}$}
    \State \textbf{(Coefficient projection with fixed-$N$ normalization)}
    \State For each $m\in\mathcal{M}_K$, set
    \[
      a_m \;\gets\; \frac{1}{N}\sum_{i\in\mathcal{I}} \psi_m\!\big(\mathbf{X}^i\big).
    \]
    \State \textbf{(Series assembly)}
    \[
      \tilde{\rho}_n(\mathbf{x}) \;\gets\; \sum_{m\in\mathcal{M}_K} a_m\,\psi_m(\mathbf{x}).
    \]
    \State \textbf{(Optional gradients for drift/evaluation)}
    \If{gradients are required and $\{\nabla\psi_m\}$ are available}
      \[
        \nabla\tilde{\rho}_n(\mathbf{x}) \;\gets\; \sum_{m\in\mathcal{M}_K} a_m\,\nabla\psi_m(\mathbf{x}).
      \]
    \EndIf
    \State \Return $\tilde{\rho}_n$ (and $\nabla\tilde{\rho}_n$ if computed)
\EndProcedure
\end{algorithmic}
\end{algorithm}

\section{Heuristic derivation of the limit dynamics}

We examine the consistency of the splitting scheme for the finite-$N$ system over one step $\tau$. For simplicity, we consider the numerical solution on a periodic domain $\mathbb{T}_L^2$. Given $\varepsilon > 0$, we assume that there exists a bounded domain $\Omega_L \in [-L, L]^2$ such that 
$
\int_{\Omega_L^{\rm c}} \rho(t, \mb x) \diff \mb x < \varepsilon
$ for $t\in [0, T]$. 
The rigorous analysis in $\mathbb{R}^2$ will be pursued in future work. 
%
For simple notation, let ${\rm e}^{-\mathcal{A}_{K,N} \tau}$ denote the finite-$N$ approximation of the advection--diffusion step ${\rm e}^{-\mathcal{A} \tau}$ and $\Psi_{{\mathcal{A}}_{K,N}}^{\tau}$ the Euler-Maruyama time discreitzation in Sec. \ref{sec:adevection_diffusion}, where $\mathcal{A}_{K,N}$ represents the corresponding operator approximated by the interpolation of the target density $\rho_n(\mb x)$ using the empirical measure $\mu_n^N$ in Sec. \ref{sec:interpolation}. Similarly, ${\rm e}^{-\mathcal{B}_{K,N} \tau}$ and $\Psi_{{\mathcal{B}}_{K,N}}^{\tau}$ denote the finite-$N$ approximation of the reaction step and the stochastic branching realization in Sec. \ref{sec:branching}. 

\begin{proposition}
Let $\mu^N_n= \frac{1}{N}\sum_{i\in \mathcal I^N_n}\delta_{\bX_{N,n}^i}$ and $\left\{\bx_{N,n}^i\right\}_{i=1}^{N}$ is a sample set of density $\rho_n(\mb x) \in H^s(\mathbb{T}_L^2)$ with $s>1$ and the coefficient function $\tilde{a}$, $\tilde{d}$ and $\tilde{r}$ are $C^3_b$ on the range of $\rho$, there exists $C$ such that
\begin{equation}
\mathbb{E}\left\vert \langle f,\mu_{n+1}^N\rangle
- \big\langle f, {\rm e}^{(\mathcal{A} + \mathcal{B})\tau} \mu_n^N\big\rangle\right\vert
\;\le\; C \Big(\tau^2 + \tau K^{-(s-1)} + \tau K^2 {\rm e}^{r_s t_n/2} N^{-1/2}\Big),
\label{eq:one_step_err}
\end{equation}
for all $f\in C^3_b$ and $r_s = \sup_{\rho} \tilde{r}(\rho)$. 
\end{proposition}


\begin{proof}
We decompose the one-step error as 
\[
\begin{split}
\left \vert \langle f,\mu_{n+1}^N\rangle
- \left\langle f, {\rm e}^{(\mathcal{A} + \mathcal{B})\tau} \mu_n^N\right\rangle \right \vert 
&\le   \left \vert \langle f,\Psi_{{\mathcal{B}}_{K,N}}^{\tau} \circ \Psi_{{\mathcal{A}}_{K,N}}^{\tau} \mu_{n}^N\rangle -  \langle f,{\rm e}^{(\mathcal{A}_{K,N} + \mathcal{B}_{K,N})\tau} \mu_{n}^N\rangle \right\vert   \\
&+ \left \vert \langle f, {\rm e}^{(\mathcal{A}_{K,N} + \mathcal{B}_{K,N})\tau} \mu_{n}^N\rangle - \big\langle f, {\rm e}^{(\mathcal{A} + \mathcal{B})\tau} \mu_n^N\big\rangle \right \vert .
\end{split}
\] 

The error of the Euler-Maruyama discretization of $\Psi_{{\mathcal{A}}_{K,N}}^{\tau}$ by Eq. \eqref{eq:euler_maruyama} is $O(\tau^2)$ in expectation. Furthermore, by Proposition \ref{prop:stochastic_branch}, the stochastic branch realization is exact in expectation, i.e.,  
$
\mathbb{E}\big[\langle f, \Psi_{{\mathcal{B}}_{K,N}}^{\tau} \mu_{n}^{\ast}\rangle\big] =  \langle f, \Phi_{{\mathcal{B}}_{K,N}}^{\tau} \mu_{n}^{\ast}\rangle 
$. Therefore,  the first order Lie-Trotter discretization error is bounded by
\begin{equation}
\mathbb{E} \left \vert \langle f,\Psi_{{\mathcal{B}}_{K,N}}^{\tau} \circ \Psi_{{\mathcal{A}}_{K,N}}^{\tau} \mu_{n}^N\rangle -  \langle f,{\rm e}^{(\mathcal{A}_{K,N} + \mathcal{B}_{K,N})\tau} \mu_{n}^N\rangle \right \vert < C_f \tau^2.
\label{eq:trotter_error}
\end{equation}

The second error term arises from the interpolation approximation of $\rho_n(\mb x)$ by the empirical measure $\mu_n^N$, i.e., 
\begin{equation}
\begin{split}
&\Big|\langle f, {\rm e}^{(\mathcal{A}_{K,N} + \mathcal{B}_{K,N})\tau} \mu_{n}^N\rangle - \big\langle f, {\rm e}^{(\mathcal{A} + \mathcal{B})\tau} \mu_n^N\big\rangle\Big| \\
\le &\ C\,\tau\left(\Vert \tilde{a}_{K,N}-\tilde{a}\Vert_{L_\infty}+\Vert\tilde{d}_{K,N}-\tilde{d}\Vert_{L_\infty}+\Vert \tilde{r}_{K,N}-\tilde{r}\Vert_{L_\infty}\right) \\
\le & \ C\, \tau \Vert\rho_n^{N,K}-\rho_n\Vert_{L_\infty},
\end{split}    
\end{equation}
where $\rho_n^{N,K}$ is the interpolation from the empirical measure $\mu_n^N$ using $2K+1$ Fourier modes. 
%
The interpolation error can be decomposed into the empirical-measure sampling and the spectral truncation error, i.e., 
\[
\Vert\rho_n^{N,K}-\rho_n\Vert_{L_\infty} \le \Vert\rho_n^{N,K}-\rho_n^{K}\Vert_{L_\infty} + \Vert\rho_n^{K}-\rho_n\Vert_{L_\infty},
\]
where $\rho_n^{K}$ represents the $K$-mode spectral representation of $\rho_n$. By Eq. \eqref{eq:Fourier_mode}, the error of the individual Fourier coefficients is bounded by 
$\vert \mathcal{I}_{n}\vert/N^2$. 
% \[
% \mathbb{E}[(\hat{\rho}^{N}(k) - \hat{\rho}(k))^2] \le \vert \mathcal{I}_{n}\vert/N^2.
% \]
Therefore, we have
\[
\mathbb{E}\Vert\rho^{N,K}_n - \rho^{K}_n)\Vert_{L_2}^2 \le C K^2 \vert \mathcal{I}_{n}\vert/N^2
\]
and
\begin{equation}
\begin{split}
\mathbb{E} \Vert\rho_n^{N,K}-\rho_n^K\Vert_{L_\infty} &\le 
CK \mathbb{E}\Vert\rho^{N,K}_n - \rho^{K}_n\Vert_{L_2} \\
&\le 
CK\left(\mathbb{E}\Vert\rho^{N,K}_n - \rho^{K}_n\Vert_{L_2}^2\right)^{1/2} \\
&\le CK^2 \vert \mathcal{I}_{n}\vert^{1/2}/N \\
&= C K^2 {\rm e}^{r_s t_n/2}/\sqrt{N}.
\end{split}
\label{eq:MC_error}
\end{equation}
The spectral truncation error is bounded by 
\begin{equation}
\begin{split}
\Vert\rho_n^{K}-\rho_n\Vert_{L_\infty} \le C K^{-(s-1)}.   
\end{split}
\label{eq:spectral_error}
\end{equation}
By using Eqs. \eqref{eq:trotter_error}, \eqref{eq:MC_error} and \eqref{eq:spectral_error}, we have Eq. \eqref{eq:one_step_err}. 
\end{proof}








\section{Error Estimations}
In this section, we analyze the error between the numerical and exact solutions for Equation (\ref{equ:reaction-diffusion}). First, we introduce the linearized PDE 
\begin{equation}\label{equ:linearlized}
    \partial \hat u  = \nabla \cdot \left( \hat{a}(t,x) \hat u + \nabla(\hat d(t,x)\hat u)\right) + \hat r(t,x) \hat u, 
\end{equation}
where $\hat a(t,x) = a (\mathcal P_K u(t,x))$, $\hat d(t,x) = d (\mathcal P_K u(t,x))$ and $\hat r(t,x) = r (\mathcal P_K u(t,x))$, and $P_K$ means the projection operator on the Fourier space upto smallest $K$ mode.
\begin{definition}[Weak formulation]
Let $\Omega\subset\mathbb{R}^2$ be a periodic domain and $T>0$.  
Assume $\hat a,\hat d,\hat r \in L^\infty(\Omega_T)$ with $\hat d(t,x)\ge d_0>0$.  
Given initial data $\hat u_0 \in L^2(\Omega)$, we call 
\[
\hat u \in L^2(0,T;H^1(\Omega)) \cap C([0,T];L^2(\Omega))
\]
a weak solution of \eqref{equ:linearlized} if for all test functions $f \in H^1(\Omega)$ and for almost every $t\in (0,T)$,
\begin{equation}\label{eq:weak-evans}
\begin{aligned}
    &\left\langle\partial_t \hat u(t), f\right\rangle+ \int_\Omega \big( \hat a(t,x)\,\hat u (x)\cdot \nabla f(x) 
+ \nabla (\hat d(t,x) \hat u (x))\cdot \nabla f(x) \big)\,dx \\
=& \left\langle\hat r(t,\cdot)\,\hat u(t,\cdot), f\right\rangle,
\end{aligned}
\end{equation}
where $\langle\cdot,\cdot\rangle$ denotes the $L^2(\Omega)$ inner product.
The initial condition is understood in the sense
\[
(\hat u(0), f)_{L^2} = (\hat u_0, f)_{L^2}, \quad \forall f\in L^2(\Omega).
\]
\end{definition}


\begin{lemma}[Well-posedness of the linearized PDE]\label{lemma:well_posedness}
Assume that $\hat a(t,x), \hat d(t,x)$,  and  $ \hat r(t,x) \in L^\infty(\Omega_T)$ with 
$\hat d(t,x) \ge d_0 > 0$ for some constant $d_0$. 
Then the linear parabolic PDE \eqref{equ:linearlized}
admits a unique weak solution 
\[
\hat u \in L^2\big(0,T; H^1(\Omega)\big) \cap C\big([0,T]; L^2(\Omega)\big),
\]
satisfying 
for any initial data $\hat u(0,\cdot) \in L^2(\Omega)$.
\end{lemma}
The proof can be found in standard PDE textbooks such as ~\cite{evans2022partial}.
The continuous limit of our algorithm is described by a stochastic differential equation
\[
\intd \bX^i_N(t) = \hat a(t,\bX^i_N(t))\intd t  + \hat d (t,\bX^i_N(t)) \intd \bW^i(t).
\]
Also, any particle $i\in \mathcal I(N,t)$ at position $\bX^i_N(t)$ and time $t$ will induce discontinuous changes to the population by giving birth a particle $i^*\in \mathcal I(N,t)$ at $\bX^i_N$ with intensity $\hat r(x)$ if $\hat r(x)>0$ and die with intensity $-\hat r(x)$ of $\hat r(x)<0$. It induces the following two point process for $\sigma\in \mathbb R^+$ 
\[
\begin{aligned}
   \mathcal R^{i^*,birth}_N( \sigma) = Q^{i,birth}_N \left(\int_{0}^{ \sigma} \mathbf 1 _{S(N,\tau)}(i)\hat r(t,\bX^i_N(\tau))\intd \tau\right),\\
    \mathcal R^{i^*,death}_N( \sigma) = Q^{i,death}_N \left(-\int_{0}^{\sigma} \mathbf 1 _{S(N,\tau)}(i)\hat r(t,\bX^i_N(\tau))\intd \tau\right), 
\end{aligned}
\]
where $Q^{i,birth}_N$, $Q^{i,death}_N$ are independent standard Poisson processes. 
\begin{proposition}
    Assume that $\hat a(t,x), \hat d(t,x)$,  and  $ \hat r(t,x) \in L^\infty(\Omega_T)$, therfore we have for $f \in H^2(\Omega) \subseteq H^1(\Omega)$,
    \[
    \mathbb E \left[ \langle f,\mu^N(t,\cdot)-\hat u(t,\cdot) \rangle \right] = 0,
    \] and 
    \[
    \mathbb E \left[ \langle f,\mu^N(t,\cdot)-\hat u(t,\cdot) \rangle^2 \right] = \frac{1}{N},
    \]
    where $\mu^N$ is the empirical measure valued stochastic process before.
\end{proposition}
\begin{proof}
    Using It\^o's formula, we have 
\begin{equation*}
    \begin{aligned}
        \langle f,\mu^N(t,\cdot)\rangle = &\left\langle f(\tau,\cdot), \mu^N(0,\cdot)\right\rangle - \int_0^t\left\langle \hat a(\tau,\cdot)\cdot \nabla f(\tau,\cdot) ,\mu^N(\tau,\cdot)\right\rangle\intd \tau \\
        &+ \int_0^t\left\langle \hat d(\tau, \cdot) \Delta f(\tau,\cdot),\mu^N(\tau,\cdot)\right\rangle \intd \tau\\
        &+ \int_0^t\left\langle  \frac{\partial}{\partial \tau} f(\tau,\cdot),\mu^N(\tau,\cdot)\right\rangle \intd\tau\\
        &+ \frac{1}{N} \int_0^t \sum_{i\in \mathcal I (N,\tau)} \sqrt{2 \hat d(\tau,\bX_N^i(\tau))} \nabla f(\tau,\bX_N^i(\tau))\cdot \intd \bW_N^i(\tau)\\
        &+  \frac{1}{N} \int_0^t \sum_{i\in \mathcal I (N,\tau), \hat r(\tau, \bX^i_N(\tau))>0} f(\tau,\bX_N^i(\tau))\mathcal R_N^{i^*,birth}(\intd \tau)\\
        &-  \frac{1}{N} \int_0^t \sum_{i\in \mathcal I (N,\tau), r(\tau, \bX^i_N(\tau))<0} f(\tau,\bX_N^i(\tau)) \mathcal R_N^{i^*,death}(\intd \tau)
    \end{aligned}
\end{equation*}
One defines 
\begin{equation*}
    \begin{aligned}
        \mathcal M^1(t,f) =& \frac{1}{N} \int_0^t \sum_{i\in \mathcal I(N,\tau)} \sqrt{2 \hat d(\tau,\bX_N^i(\tau))} \nabla f(\bX_N^i(\tau))\cdot \intd \bW_N^i(\tau)\\
        \mathcal M^2(t,f) =& \frac{1}{N} \int_0^t \sum_{i\in \mathcal I(N,\tau), r(\bX^i_N(\tau))>0} f(\bX_N^i(\tau))\left(\mathcal R_N^{i^*,birth}(\intd \tau) -r(\bX_N^i(\tau))\intd \tau\right)\\
        \mathcal M^3(t,f) =& \frac{1}{N} \int_0^t \sum_{i\in \mathcal I(N,\tau), r(\bX^i_N(\tau))<0} f(\bX_N^i(\tau)) \left(\mathcal R_N^{i^*,death}(\intd \tau)-r(\bX_N^i(\tau))\intd \tau\right).
    \end{aligned}
\end{equation*}
We notice that 
\begin{equation}
\begin{aligned}
    \mathbb E\left[ \left(\mathcal M^1(t,f) \right)\right] &= 0
\end{aligned}
\end{equation}
and
\begin{equation}
\begin{aligned}
    \mathbb E\left[ \left(\mathcal M^1(t,f) \right)^2\right] &= \frac{1}{N^2} \int_0^t \sum_{i\in \mathcal I(N,\tau)}  \left|\hat d(t,\bX^i_N(\tau))\nabla f(\bX_N^i(\tau))\right|^2 \intd \tau \\
    & \leq \frac{C_{f',\sigma}t}{N},
\end{aligned}
\end{equation}
where $C_{f',\sigma}$ is a constant only dependent on the gradient of $f$ and $\sigma$. Similarily, we have
\begin{equation}
\begin{aligned}
    \mathbb E\left[ \mathcal M^2(t,f) \right] &= 0
\end{aligned}
\end{equation}
and 
\begin{equation*}
    \begin{aligned}
       \mathbb E\left[ \left(\mathcal M^2(t,f) \right)^2\right] &= \frac{1}{N^2}  \mathbb E\left[ \int_0^t \sum_{i\in \mathcal I(N,\tau), r(\bX^i_N(\tau))>0}  f(\bX_N^i(\tau))\mathcal R_N^{i^*,birth}(\intd \tau) \cdot\right.\\
    & \qquad \left.\int_0^{t'} \sum_{j\in \mathcal I(N,\tau), r(\bX^j_N(\tau'))>0}  f(\bX_N^j(\tau'))\mathcal R_N^{j^*,birth}(\intd \tau')\right]\\
       & - \frac{2}{N^2} \mathbb E\left[\int_0^t \sum_{i\in \mathcal I(N,\tau), r(\bX^i_N(\tau))>0}  f(\bX_N^i(\tau))\mathcal R_N^{i^*,birth}(\intd \tau)\right.\\
    & \qquad \left.\int_0^{t'} \sum_{j\in \mathcal I(N,\tau), r(\bX^j_N(\tau'))>0}  f(\bX_N^j(\tau'))r(\bX_N^j(\tau'))\intd \tau'\right]\\
       & + \frac{1}{N^2} \left(\int_0^{t} \sum_{i\in \mathcal I(N,\tau), r(\bX^i_N(\tau))>0}  f(\bX_N^i(\tau))r(\bX_N^i(\tau))\intd \tau\right)^2\\
       & = \frac{1}{N^2}I_1-\frac{2}{N^2}I_2+\frac{1}{N^2}I_3.
    \end{aligned}
\end{equation*}
For $I_1$, we have 
\begin{equation*}
    \begin{aligned}
        I_1 =& \sum_{i\in \mathcal I(N,\tau), r(\bX^i_N(\tau))>0}\int_0^{t}f^2(\bX_N^i(\tau))r(\bX_N^i(\tau)) \intd \tau\\
        &+\left(\int_0^{t} \sum_{i\in \mathcal I(N,\tau), r(\bX^i_N(\tau))>0}  f(\bX_N^i(\tau))r(\bX_N^i(\tau))\intd \tau\right)^2.
    \end{aligned}
\end{equation*}
For $I_2$, we have
\begin{equation*}
    \begin{aligned}
        I_2 &= \left(\int_0^{t} \sum_{i\in \mathcal I(N,\tau), r(\bX^i_N(\tau))>0}  f(\bX_N^i(\tau))r(\bX_N^i(\tau))\intd \tau\right)^2.
    \end{aligned}
\end{equation*}
Therefore, we have
$$
\begin{aligned}
\mathbb E\left[ \left(\mathcal M^2(t,f) \right)^2\right] &= \frac{1}{N^2}\sum_{i\in \mathcal I(N,\tau), r(\bX^i_N(\tau))>0}\int_0^{t}f^2(\bX_N^i(\tau))r(\bX_N^i(\tau)) \intd \tau\\
    & = \frac{t C_{f,r}}{N},
\end{aligned}
$$
where $C_{f,r}$ is a constant dependent on $f$ and $r$. Similarly, we can show that 
$$
\begin{aligned}
\mathbb E\left[ \mathcal M^3(t,f) \right] 
    & = 0,
\end{aligned}
$$
and 
$$
\begin{aligned}
\mathbb E\left[ \left(\mathcal M^3(t,f) \right)^2\right] 
    & = \frac{t C_{f,r}}{N}.
\end{aligned}
$$
% As a result, we have
% \begin{equation*}
%     \begin{aligned}
%          \langle f(t,\cdot),\mu^N(t,\cdot)\rangle = &\langle f(0,\cdot), \mu^N(0,\cdot)\rangle \\
%          & + \int_0^t\langle -\hat a(\cdot)\cdot \nabla f (\tau,\cdot) + \hat d(\tau, \cdot) \Delta f(\tau,\cdot) + r(\cdot) ,\mu^N(\tau,\cdot)\rangle \intd \tau\\
%          & + \int_0^t\langle \frac{\partial}{\partial \tau} f(\tau,\cdot),\mu^N(\tau,\cdot)\rangle \intd \tau \\
%          & + \mathcal M^1+\mathcal M^2 + \mathcal M^3.
%     \end{aligned}
% \end{equation*}
Notice that 
\begin{equation*}
    \begin{aligned}
         \langle f,\mu(t,\cdot)-\hat u(t,\cdot) \rangle = &\langle f, \mu^N(0,\cdot)-\hat u(0,\cdot)\rangle \\
         &  +\int_0^t\langle -\hat a(\tau ,\cdot )\cdot \nabla f  + \hat d(\tau, \cdot) \Delta f+r(\tau,\cdot),\mu(\tau,\cdot)-\hat u(\tau,\cdot)\rangle\intd \tau\\
         & + \int_0^t\langle f ,\frac{\partial}{\partial \tau} \mu(\tau,\cdot)-\frac{\partial}{\partial \tau} \hat u(\tau,\cdot)\rangle\intd \tau\\
         & + \mathcal M^1+\mathcal M^2 + \mathcal M^3.
    \end{aligned}
\end{equation*}
Therefore, by taking the expectation, we have 
\begin{equation}
\begin{aligned}
        \mathbb E \left[\left\langle f,\mu(t,\cdot) \right\rangle\right] = &
          \mathbb E \left[\langle f, \mu^N(0,\cdot) \rangle\right] \\
         &  +\int_0^t\mathbb E \left[ \langle -\hat a(\tau,\cdot)\cdot \nabla f  + \hat d(\tau,\cdot) \Delta f +\hat r(\tau,\cdot),\mu(\tau,\cdot)\rangle\right]\intd \tau\\
         & + \int_0^t\mathbb E \left[ \langle f ,\frac{\partial}{\partial \tau} \mu(\tau,\cdot)\rangle\right]\intd \tau\\
\end{aligned}
\end{equation}
Based on the uniqueness of the weak solution of the \eqref{equ:linearlized} in Lemma \ref{lemma:well_posedness}, we  have 
\begin{equation}
    \mathbb E \left[\left\langle f,\mu(t,\cdot)-\hat u(t,\cdot) \right\rangle\right] = 0.
\end{equation}
For the second moment, the quadratic variation of each martingale term gives
\[
\mathbb{E}\!\left[ M_1(t,f)^2 \right] \;\le\; \frac{C_{f,\sigma}\,t}{N}, 
\quad 
\mathbb{E}\!\left[ M_2(t,f)^2 \right] \;\le\; \frac{C_{f,r}\,t}{N}, 
\quad 
\mathbb{E}\!\left[ M_3(t,f)^2 \right] \;\le\; \frac{C_{f,r}\,t}{N},
\]
where $C_{f,\sigma}$ and $C_{f,r}$ are constants depending on the test function 
$f$, the bounded diffusion coefficient, and the reaction coefficient. 
Cross-variations vanish by independence. Therefore,
\[
\mathbb{E}\!\left[\langle f(t), \mu^N(t) - \hat u(t) \rangle^2\right]
= \mathbb{E}\!\left[ M_1(t,f)^2 + M_2(t,f)^2 + M_3(t,f)^2 \right] 
\;\le\; \frac{C\,t}{N},
\]
which establishes the claimed $\mathcal{O}(1/N)$ rate.
\end{proof}
% \newpage

% \begin{proof}
%     Using It\^o's formula, we have 
% \begin{equation*}
%     \begin{aligned}
%         \langle f(t,\cdot),\hat{\mu}^N(t,\cdot)\rangle = &\left\langle f(0,\cdot), \hat{\mu}^N(0,\cdot)\right\rangle - \int_0^t\left\langle a\left(\left[\mathcal P_K \hat \mu^N\right](\cdot) \right)\cdot \nabla f (\tau,\cdot),\hat{\mu}^N(\tau,\cdot)\right\rangle\intd \tau \\
%         &+ \int_0^t\left\langle d\left(\left[\mathcal P_K \hat \mu^N\right](\cdot) \right) \Delta f(\tau,\cdot),\hat \mu^N(\tau,\cdot)\right\rangle \intd \tau\\
%         &+ \int_0^t\left\langle  \frac{\partial}{\partial \tau} f(\tau,\cdot),\hat \mu^N(\tau,\cdot)\right\rangle \intd\tau\\
%         &+ \frac{1}{N} \int_0^t \sum_{i\in \mathcal I (N,\tau)} \sqrt{2 d\left(\left[\mathcal P_K \hat \mu^N\right]\left(\bX_N^i(\tau)\right) \right)} \nabla f(\tau,\bX_N^i(\tau))\cdot \intd \bW_N^i(\tau)\\
%         &+  \frac{1}{N} \int_0^t \sum_{i\in \mathcal I (N,\tau), r\left(\left[\mathcal P_K \hat \mu^N\right]\left(\bX_N^i(\tau)\right) \right)>0} f(\tau,\bX_N^i(\tau))\mathcal R_N^{i^*,birth}(\intd \tau)\\
%         &-  \frac{1}{N} \int_0^t \sum_{i\in \mathcal I (N,\tau), r\left(\left[\mathcal P_K \hat \mu^N\right]\left(\bX_N^i(\tau)\right) \right)<0} f(\tau,\bX_N^i(\tau)) \mathcal R_N^{i^*,death}(\intd \tau)
%     \end{aligned}
% \end{equation*}
% Subtract the exact solution, gives
% \begin{equation*}
%     \begin{aligned}
%         &\langle f(t,\cdot),\hat{\mu}^N(t,\cdot)-u(t,\cdot)\rangle \\= &\left\langle f(0,\cdot), \hat{\mu}^N(0,\cdot)-u(0,\cdot)\right\rangle \\
%         & - \int_0^t\left\langle a\left(\left[\mathcal P_K \hat \mu^N\right](\cdot) \right)\cdot \nabla f (\tau,\cdot),\hat{\mu}^N(\tau,\cdot)\right\rangle - \left\langle a(u(\tau,\cdot))\cdot \nabla f(\tau,\cdot) ,u(\tau,\cdot) \right\rangle \intd \tau \\
%         &+ \int_0^t\left\langle d\left(\left[\mathcal P_K \hat \mu^N\right](\cdot) \right) \Delta f(\tau,\cdot),\hat \mu^N(\tau,\cdot)\right\rangle - 
%         \left\langle d\left(u(\tau,\cdot)  \right) \Delta f(\tau,\cdot),u(\tau,\cdot)\right\rangle
%         \intd \tau\\
%         &+ \int_0^t\left\langle  \frac{\partial}{\partial \tau} f(\tau,\cdot),\hat \mu^N(\tau,\cdot) -  u(\tau,\cdot)\right\rangle \intd\tau\\
%         &+ \frac{1}{N} \int_0^t \sum_{i\in \mathcal I (N,\tau)} \sqrt{2 d\left(\left[\mathcal P_K \hat \mu^N\right]\left(\bX_N^i(\tau)\right) \right)} \nabla f(\tau,\bX_N^i(\tau))\cdot \intd \bW_N^i(\tau)\\
%         &+  \frac{1}{N} \int_0^t \sum_{i\in \mathcal I (N,\tau), r\left(\left[\mathcal P_K \hat \mu^N\right]\left(\bX_N^i(\tau)\right) \right)>0} f(\tau,\bX_N^i(\tau))\mathcal R_N^{i^*,birth}(\intd \tau) \\
%         & - \int_0^t  \int_{r(u(\tau,\bx))>0} f(\tau,\cdot), r(u(\tau,\cdot))u(\tau,\cdot)\intd \bx \intd \tau \\
%         &-  \frac{1}{N} \int_0^t \sum_{i\in \mathcal I (N,\tau), r\left(\left[\mathcal P_K \hat \mu^N\right]\left(\bX_N^i(\tau)\right) \right)<0} f(\tau,\bX_N^i(\tau)) \mathcal R_N^{i^*,death}(\intd \tau)\\
%         & - \int_0^t  \int_{r(u(\tau,\bx))<0} f(\tau,\cdot), r(u(\tau,\cdot))u(\tau,\cdot)\intd \bx \intd \tau \\
%         & = \mathcal  E_0 + \mathcal E_a + \mathcal E_d + \mathcal E_f + \mathcal N_w +\mathcal N_p +\mathcal N_n.
%     \end{aligned}
% \end{equation*}
% For the $\mathcal E_a$, we have that
% \[
% \begin{aligned}
% \mathcal E_a = &  \int_{0}^t \left\langle a\left(\left[\mathcal P_K \hat \mu^N\right](\cdot) \right)\cdot \nabla f (\tau,\cdot),\hat{\mu}^N(\tau,\cdot)\right\rangle - \left\langle a(u(\tau,\cdot))\cdot \nabla f(\tau,\cdot) ,u(\tau,\cdot) \right\rangle \intd \tau\\
%      = & \int_{0}^t \left\langle a\left(\left[\mathcal P_K \hat \mu^N\right](\cdot) \right)\cdot \nabla f (\tau,\cdot),\hat{\mu}^N(\tau,\cdot)  - u(\tau,\cdot) \right\rangle  \intd \tau
%     \\ & + \int_{0}^t  \left\langle \left(a\left(\left[\mathcal P_K \hat \mu^N\right](\cdot) \right) -a(u(\tau,\cdot))\right)\cdot \nabla f(\tau,\cdot) , u(\tau,\cdot) \right\rangle \intd \tau
% \end{aligned}
% \]

% \end{proof}

% \newpage
Next, we estimate the error between the linearized PDE \eqref{equ:linearlized} and the original PDE \eqref{equ:reaction-diffusion}.\begin{proposition}
Assume that the coefficients $a(u), d(u), r(u)$ are $C^1$ functions with bounded derivatives, and that $d(u)\ge d_0>0$ uniformly. 
Suppose $u$ is a sufficiently smooth solution of the nonlinear PDE~\eqref{equ:reaction-diffusion} on $\Omega$ with periodic  boundary conditions, and $\hat u$ is the solution of the linearized PDE~\eqref{equ:linearlized}. 
If $u(\cdot,t)\in H^m(\Omega)$ for some $m>1$, then the error $e=u-\hat u$ satisfies
\[
\|e(t)\|_{L^2}^2 \;\le\; C_T\Big(\|e(0)\|_{L^2}^2 + K^{2(1-m)}\|u\|_{L^2(0,t;H^m)}^2\Big),
\]
for some constant $C_T>0$ depending only on $T$ and the coefficients.
\end{proposition}


\begin{proof}
    Let $\eta(t) = \frac{1}{2}\|e(t,\cdot)\|^2_{L^2(\Omega)}  = \frac{1}{2}\|u(t,\cdot)-\hat u(t,\cdot)\|^2_{L^2(\Omega)} $, then we have 
    \begin{equation}
        \begin{aligned}
    \frac{\partial}{\partial t} \eta(t)
        = & \int_\Omega (u-\hat{u})\nabla \cdot (a(u)u - a(\mathcal P_K u )\hat u )\intd \bx \\
        & + \int_\Omega (u-\hat{u})  \Delta (d(u)u - d(\mathcal P_K u )\hat u  )\intd \bx \\
        & + \int_\Omega (u-\hat{u})  (r(u)u - r(\mathcal P_k  u )\hat u) \intd \bx  \\
        = &  I_1 + I_2 + I_3
        \end{aligned}
    \end{equation}
Notice that for $I_1$, we have
\[
\begin{aligned}
    I_1 =& \int_\Omega (u-\hat{u})\nabla \cdot (a(u)u - a(u)\hat u )\intd \bx \\
    &+ \int_\Omega (u-\hat{u})\nabla \cdot (a(u)\hat u - a(\mathcal P_K u )\hat u )\intd \bx = I_{1,1}+ I_{1,2}.\\
\end{aligned}
\]
For $I_{1,1}$, we have
\[
\begin{aligned}
    I_{1,1} & = \int_\Omega (u-\hat{u})^2\nabla \cdot a(u)\intd \bx \\
    & +  \int_\Omega (u-\hat{u})a(u)\cdot\nabla (u-\hat{u}) \intd \bx \\
    = &  \int_\Omega (u-\hat{u})^2\nabla \cdot a(u)\intd \bx \\
    & +  \frac{1}{2}\int_\Omega a(u)\cdot\nabla (u-\hat{u})^2 \intd \bx \\
   = & \frac{1}{2}\int_\Omega e^2\nabla \cdot a(u)\intd \bx
    \end{aligned}
\]
Therefore we have\[
|I_{1,1}|  
    \leq \frac{1}{2}\|a'\|_{\infty}\|\nabla u \|_{L^\infty}\|e\|_{L^2} = \|a'\|_{\infty}\|\nabla u \|_{L^\infty} \eta.
\]
Also, for $I_{1,2}$, we have 
\[
\begin{aligned}
    I_{1,2} = & - \int_\Omega \nabla (u-\hat{u}) \cdot (a(u)\hat u - a(\mathcal P_K u )\hat u )\intd \bx
\end{aligned}
\]
Therefore, we using Cauchy–Schwarz–Young to get 
\[
\begin{aligned}
|I_{1,2}| \leq& \|\hat u \|_{L^\infty} \|\nabla e \|_{L^2}\|a(u) - a(\mathcal P_K u )\|_{L^2}\\
\leq & \|\hat u \|_{L^\infty} \|a'\|_\infty \|\nabla e \|_{L^2}\|u -\mathcal P_K u \|_{L^2}\\
\leq &  \epsilon \|\nabla e \|_{L^2}\| + C_\epsilon \|u -\mathcal P_K u \|_{L^2}.
\end{aligned}
\]
For $I_2$, we divide the diffusion term as
\[
\nabla d(u)u = (d(u) + d'(u)u)\nabla u.
\]
Therefore, 
\[
\begin{aligned}
    I_2 = &  -\int \nabla e \cdot \nabla  (d(u)u - d(\mathcal P_K u )\hat u  )\intd \bx\\
     = & -\int \nabla e \cdot \left( (d(u)+d'(u)u)\nabla u - d(\mathcal P_K u) \nabla \hat u - \hat u d'(\mathcal P_K u) \nabla (\mathcal P_K u)  \right)\intd \bx \\
     = &-\int   d(\mathcal P_K u) \nabla e\cdot \nabla ( u - \hat u ) -  \int (d(u) - d(\mathcal P_K u) ) \nabla e \cdot  \nabla u \\
     & - \int \nabla e \cdot \left(  d'(u)u\nabla u- \hat u d'(\mathcal P_K u) \nabla (\mathcal P_K u)\right)\intd \bx  \\
    = & -\int   d(\mathcal P_K u) \nabla e\cdot \nabla ( u - \hat u ) -  \int (d(u) - d(\mathcal P_K u) ) \nabla e \cdot  \nabla u \\
     & -  \int \nabla e \cdot \left(  d'(u)u\nabla u- \hat u d'(u) \nabla u \right)\intd \bx\\
     & -\int \nabla e \cdot \left( \hat u d'(u) \nabla u- \hat u d'(\mathcal P_K u) \nabla u \right)\intd \bx\\
     & -  \int \nabla e \cdot \left(   \hat u d'(\mathcal P_K u) \nabla u - \hat u d'(\mathcal P_K u) \nabla (\mathcal P_K u)\right)\intd \bx\\
     =& I_{2,1}+I_{2,2}+I_{2,3}+I_{2,4}+I_{2,5}. 
\end{aligned}
\]
Notice that 
\[
\begin{aligned}
    I_{2,1} &\leq -d_0\|\nabla e \|_{L^2},\\
I_{2,2} &\leq \epsilon \|\nabla e \|_{L^2} +C_\epsilon |d'|_\infty \|u -\mathcal P_K u \|_{L^2} \|\nabla u\|_{L^\infty},\\
I_{2,3} &\leq  \epsilon \|\nabla e \|_{L^2} + C_\epsilon \|e \|_{L^2} |d'|_\infty\|\nabla u\|_{L^\infty},\\
I_{2,4} &\leq  \epsilon \|\nabla e \|_{L^2} + C_\epsilon\|\hat u \|_{L^\infty} |d''|_\infty\|u -\mathcal P_K u \|_{L^2}\|\nabla u\|_{L^\infty},\\
I_{2,5} &\leq  \epsilon \|\nabla e \|_{L^2} + C_\epsilon \|\nabla u -\nabla \mathcal P_K u  \|_{L^2}\|\hat u \|_{L^\infty}\|d'(\mathcal P_K u) \|_{L^\infty}.
\end{aligned}
\]
In total, we will have
\[
I_2 \leq -(d_0-4\epsilon)\|\nabla e \|_{L^2}  + C_1 \|e\|_{L^2} + C_2 \| u - \mathcal P_K u  \|_{L^2}  + C_3 \| \nabla u - \nabla \mathcal P_K u  \|_{L^2}.
\]
Finally, for $I_3$, we have
\[
\begin{aligned}
I_3 =& \int e (r(u)u - r(\mathcal P_K u )\hat u \intd \bx\\
    =&  \int e^2 r(u) \intd \bx 
    + \int  e \hat{u} (r(u)- r(\mathcal P_K u )\intd \bx 
    \\
    \leq & |r|_\infty \|e\|_{L^2} + |r'|_\infty \|\hat u \|_{L^\infty}\|u - P_K u \|_{L^2}\|e\|_{L^2}\\
   \leq & C_3 \|e\|_{L^2}+C_4 \|u - P_K u \|_{L^2}\|e\|_{L^2} \\
\end{aligned}
\]
By choose $\epsilon$ small enough compared with $d_0$, we can have
\[
\frac{\partial }{\partial t }\eta \leq C_\eta \eta(t) + C_{\mathcal P } \left(\|u - P_K u \|^2_{L^2} + \|u - P_K u \|^2_{H^1} \right)
\]
Applying the Gr\"onwall inequility, we have 
\[
\|u(t) - \hat u (t)\|^2_{L^2} \leq \exp( C_\eta t )\left(\|u(0) - \hat u (0)\|^2_{L^2} +  C_{\mathcal P }\int _0^t \|u(\tau) - P_K u(\tau) \|^2_{L^2} + \|u(\tau) - P_K u(\tau) \|^2_{H^1} \intd \tau \right)
\]
Assuming that $u(t,\cdot)\in H^m(\Omega)$, $m>0$, we will have 
\[
\|u(t) - \hat u (t)\|^2_{L^2} \leq C_T\left(\|u(0) - \hat u (0)\|^2_{L^2} +  K^{2-2m}\|u\|_{L^2(0,t;H^m)} \right)
\]
\end{proof}


Combined weak error from Propositions 3 and 4, we have for every $f\in H^2(\Omega)$ and all $t\in[0,T]$,
\[
\mathbb{E}\,\big[\langle f,\mu_N(t)-u(t)\rangle^2\big]
\;\le\;
C_t\Big(
\frac{1}{N}\,\|f\|_{H^2}^2
\;+\;
\|f\|_{H^2}^2\big(\|u(0)-\hat u(0)\|_{L^2}^2
+ K^{2(1-m)}\|u\|_{L^2(0,t;H^m)}^2\big)
\Big).
\]
In particular (taking the supremum over $\|f\|_{H^2}\le1$), one obtains the $H^{-2}$-error estimate
\[
\mathbb{E}\,\|\mu_N(t)-u(t)\|_{H^{-2}}^2
\;\le\;
C_t\Big(
\frac{1}{N}
+ \|u(0)-\hat u(0)\|_{L^2}^2
+ K^{2(1-m)}\|u\|_{L^2(0,t;H^m)}^2
\Big).
\]




The numerical details to implement this can be seen in Algorithm \ref{alg:linear}.
\begin{algorithm}
    \caption{Stochastic Branching for linear scalar reaction--diffusion--advection equation} \label{alg:linear}
    \begin{algorithmic}
        \Require Initial condition of PDE $u_0(\bx)$, domain descriptor $\Omega$;
        \Require Fixed initial population size $N$
        \Require Time step $\tau>0$, horizon $T$; truncation $\mathcal M_K$; basis $\{\psi_m\}_{m\in\mathcal M_K}$
        \Require Driven term $b(\bx)$
        \Require Diffusion strength $\sigma(\bx)$
        \Require Reaction maps $r(\bx)$
        \Ensure  Solution set of PDE $\mathcal U = \{u_n\}_{n=0}^{T/\tau}$
        \Procedure{Stochastic\_Branching}{$u_0$,$N$,$\tau$,$\mathcal M_K$,$\{\psi_m\}$,$b$,$\sigma$,$r$}
        \State \textbf{Initialization}
        \State $n\gets 0$
        \State Precomputed normalization constants  $Z\gets \int_\Omega u_0(\bx)\intd \bx$  (either numerically or analytically).
        \State Generate $N$ initial particle positions set $S_0=\{\bX^i_0\}$ according to the distributions $\rho_0=\frac{u_0}{Z}$ and using the Metropolis--Hastings algorithm.
        \State \textbf{Interpolate (fixed-$N$):}
        \State $\tilde\rho_0 \gets \textsc{Interpolate}(S_0,N,\mathcal M_K,\{\psi_m\})$
        \State Append $\tilde u_0 = Z\tilde \rho_0$ into $\mathcal U$
        \While {$n\tau<T$}
            \State \textbf{SDE propagation:}
            \State $S_n^*\gets \textsc{SDE\_Propagate}(S_n,\tau,b,\sigma)$
            \State \textbf{Birth--Death:}
            \State $S_{n+1}\gets \textsc{Birth\_Death}(S_{n},\tau,r)$
            \State \textbf{refresh fields:}
            \State $\tilde\rho_{n+1} \gets \textsc{Interpolate}(S_{n+1},N,\mathcal M_K,\{\psi_m\})$
            \State Append $\tilde u_{n+1} = Z\tilde \rho_{n+1}$ into $\mathcal U$
            \State $n\gets n+1$
        \EndWhile
        \State \Return $\mathcal U$
        \EndProcedure
    \end{algorithmic}
\end{algorithm}
\section{Extension to PDE system}
\label{sec:extension}
To solve PDE systems like Equation \eqref{equ:KS}, we introduce two sets of particles $\bX^i(t)$ and $\bY^i(t)$. Therefore, empirical distribution $\rho_u(t,\intd \bx ) = \mu_u(t,\intd \bx) =  \frac{1}{N_{x,0}}\sum_{i=1}^{N_{x,n}}\delta_{\bX_n^i} $ and $\rho_v(t,\intd \bx ) = \mu_v(t,\intd \bx) =  \frac{1}{N_{y,0}}\sum_{i=1}^{N_{y,n}}\delta_{\bY_n^i}$ are used to approximate $u$ and $v$ separately, by 
\begin{equation}
    u(t,\bx) =  Z_u \rho_u(t,\bx), \quad v(t,\bx) =  Z_v \rho_v(t,\bx),
\end{equation}
where $Z_u = \int_\Omega u(0,\bx) \intd \bx$ and $Z_v = \int_\Omega v(0,\bx) \intd \bx$. Therefore, the dynamics of \eqref{equ:KS} can be written as 
% \begin{equation}
% \begin{aligned}
%     \partial_t u - \nabla \cdot ( \nabla u - \chi(u) \nabla v) &= f_u(u,v),  & x\in \mathbb R^2 ,t\in(0,T),\\
%     v_t -\Delta v & = f_v(u,v),  &x\in \mathbb R^2 , t\in(0,T),
%     \end{aligned}
% \end{equation}
\begin{equation}
    \begin{aligned}
        \partial_t \rho_u =  \mathcal A_u \rho_u + \mathcal B_u \rho_u,\\
        \partial_t \rho_v =  \mathcal A_v \rho_v + \mathcal B_v \rho_v
    \end{aligned}
\end{equation}
where $\mathcal A_u \rho = \nabla \cdot ( \nabla \rho - \frac{Z_v}{Z_u} \chi(Z_u\rho) \nabla \rho_v)$ , $\mathcal A_v \rho = \Delta \rho$, $\mathcal B_u \rho_u  = f_u(Z_u\rho_u,Z_v \rho_v)$ and $B_v \rho_v = f_v(Z_u\rho_u,Z_v \rho_v)$. Similar to before, we can approximate the solution by operator splitting. 
\begin{equation}
    \begin{aligned}
        \rho_{u,n}^*  =  \exp(\mathcal A_u\tau) \rho_{u,n}, \quad 
         \rho_{v,n}^*  =  \exp(\mathcal A_v\tau) \rho_{v,n},  \\
         \rho_{u,n+1} =  \exp(\mathcal B_u\tau) \rho_{u,n}^*,  \quad 
         \rho_{v,n+1}  =  \exp(\mathcal B_v\tau) \rho_{v,n}^*,   \\
    \end{aligned}
\end{equation}
where by matching parameters, we have $a_u(\bx;\rho_u,\rho_v) =  - \frac{Z_v}{Z_u\rho_u} \chi(Z_u\rho_u) \nabla \rho_v $ and $a_v(\bx;\rho_u,\rho_v)= 0$. For the diffusion part, we have $\sigma_u(\bx;\rho_u) = \sigma_v(\bx;\rho_v) = 1$. For birth and death part, $\tilde r_u (\bx;\rho_u,\rho_v)=  \frac{f_u(Z_u\rho_u,Z_v \rho_v)}{\rho_u}  $ and $\tilde r_v (\bx;\rho_u,\rho_v )=  \frac{f_v(Z_u\rho_u,Z_v \rho_v)}{\rho_v}  $. The detailed algorithm is shown in  Algorithm \ref{alg:ks}.

\begin{algorithm}
\caption{Stochastic Branching for keller--Segel equation}\label{alg:ks}
\begin{algorithmic}
\Require Initial PDE data $u_0(\bx),\,v_0(\bx)$; domain descriptor $\Omega$
\Require Fixed initial population sizes $N_u, N_v$
\Require Time step $\tau>0$, horizon $T$; truncation $\mathcal M_K$; basis $\{\psi_m\}_{m\in\mathcal M_K}$
% \Require Diffusion strengths $\sigma_u(\bx),\,\sigma_v(\bx)$ \hfill (constant or space-dependent)
% \Require Chemotactic coefficient: $\chi$
\Require Normalization constants $Z_u=\int_\Omega u_0\,\mathrm d\bx$, $Z_v=\int_\Omega v_0\,\mathrm d\bx$
% \Require Reaction maps $f_u,f_v:\mathbb R_{\ge 0}\times\mathbb R_{\ge 0}\to\mathbb R$
\Ensure Discrete KS fields $\mathcal U=\{\tilde u_n\}_{n=0}^{T/\tau}$, $\mathcal V=\{\tilde v_n\}_{n=0}^{T/\tau}$
% \Procedure{KS\_Stochastic\_Branching}{$u_0,v_0,N_u,N_v,\tau,\mathcal M_K,\newline \phantom{\quad\quad \textbf{procesure} \texttt{KS\_Stochastic\_Branching}(}\{\psi_m\},\sigma_u,\sigma_v,\chi,f_u,f_v$}
\Procedure{KS\_Stochastic\_Branching}{$u_0,v_0,N_u,N_v,\tau,\mathcal M_K,\{\psi_m\}$}
  \State \textbf{Initialization}
  \State $n\gets 0$
  \State Generate $N_u$,$N_v$ initial particle positions $S_{u,0}=\{\mathbf X^i_0\}$, $S_{v,0}=\{\mathbf Y^i_0\}$ via MH from $\rho_{u,0}=u_0/Z_u$.
  \State \textbf{Interpolate (fixed-$N$):}
  \State $\tilde\rho_{u,0} \gets \textsc{Interpolate}\big(S_{u,0},N_u,\mathcal M_K,\{\psi_m\}\big)$
  \State $\tilde\rho_{v,0} \gets \textsc{Interpolate}\big(S_{v,0},N_v,\mathcal M_K,\{\psi_m\}\big)$
  \State Append $\tilde u_0 = Z_u \tilde\rho_{u,0}$ to $\mathcal U$;\quad Append $\tilde v_0 = Z_v \tilde\rho_{v,0}$ to $\mathcal V$
  \While{$n\tau<T$}
    \State \textbf{SDE propagation}
    \State $b_u(\bx) \gets -\dfrac{Z_v}{Z_u\,\tilde\rho_{u,n}(\bx)}\,\chi\!\big(Z_u\tilde\rho_{u,n}(\bx)\big)\,\nabla\tilde\rho_{v,n}(\bx)$; $b_v(\bx) \gets 0$
    \State $S_{u,n,*} \gets \textsc{SDE\_Propagate}\big(S_{u,n},\tau,b_u,\sigma_u\big)$
    \State $S_{v,n,*} \gets \textsc{SDE\_Propagate}\big(S_{v,n},\tau,b_v,\sigma_v\big)$
    \State \textbf{Interpolate at starred positions (fixed-$N$)}
    \State $\tilde\rho_u^* \gets \textsc{Interpolate}\big(S_{u,n,*},N_u,\mathcal M_K,\{\psi_m\}\big)$
    \State $\tilde\rho_v^* \gets \textsc{Interpolate}\big(S_{v,n,*},N_v,\mathcal M_K,\{\psi_m\}\big)$
    \State \textbf{Birth--Death}
    \State $r_u(\bx) \gets \dfrac{f_u\!\big(Z_u\tilde\rho_u^*(\bx),\,Z_v\tilde\rho_v^*(\bx)\big)}{\tilde\rho_u^*(\bx)}$; $r_v(\bx) \gets \dfrac{f_v\!\big(Z_u\tilde\rho_u^*(\bx),\,Z_v\tilde\rho_v^*(\bx)\big)}{\tilde\rho_v^*(\bx)}$
    \State $S_{u,n+1} \gets \textsc{Birth\_Death}\big(S_{u,n,*},\tau,r_u\big)$
    \State $S_{v,n+1} \gets \textsc{Birth\_Death}\big(S_{v,n,*},\tau,r_v\big)$
    \State \textbf{Refresh fields for output (fixed-$N$)}
    \State $\tilde\rho_{u,n+1} \gets \textsc{Interpolate}\big(S_{u,n+1},N_u,\mathcal M_K,\{\psi_m\}\big)$
    \State $\tilde\rho_{v,n+1} \gets \textsc{Interpolate}\big(S_{v,n+1},N_v,\mathcal M_K,\{\psi_m\}\big)$
    \State Append $\tilde u_{n+1}=Z_u\tilde\rho_{u,n+1}$ to $\mathcal U$;\quad Append $\tilde v_{n+1}=Z_v\tilde\rho_{v,n+1}$ to $\mathcal V$
    \State $n\gets n+1$
  \EndWhile
  \State \Return $\mathcal U,\ \mathcal V$
\EndProcedure
\end{algorithmic}
\end{algorithm}



% \begin{algorithm}
% \caption{Stochastic Branching Method for the KS Equation}\label{alg:cap}
% \begin{algorithmic}
% \Require Initial population size $N_{x,0}, N_{y,0}$, time step size $\tau$, current step $n = 0$, and Fourier mode truncation $K$.
% \Require Precomputed normalization constants $Z_u$ and $Z_v$ (either numerically or analytically).
% \State Generate $N_{x,0}$ and $N_{y,0}$ initial particle positions 
% $\{\bX^i_0\}$ and $\{\bY^i_0\}$ according to the distributions 
% $\rho_u$ and $\rho_v$, respectively, using the Metropolis--Hastings algorithm.
% \State Interpolate the particle positions $\bX^i_n$ and $\bY^i_n$ onto Fourier modes up to the truncation $K$ to obtain the Fourier representations $\tilde{\rho}_u$ and $\tilde{\rho}_v$.

% \While{$n \tau < T$} \Comment{Assuming $T \mid \tau$}
%     \State Update the positions of the particles:
%     \For{$i=1$ to $N_{x,n}$}
%         \[
%         \bX^i_{n,*} = \bX^i_n - \frac{Z_v}{Z_u \tilde{\rho}_u(\bX^i_n)} \chi(Z_u \tilde{\rho}_u(\bX^i_n)) \nabla \tilde{\rho}_v(\bX^i_n) \tau + \sqrt{2\tau} \xi^{i}_{u,n}
%         \]
%         where $\xi^{i}_{u,n} \sim \mathcal{N}(0, \mathbb{I})$.
%     \EndFor
%     \For{$i=1$ to $N_{y,n}$}
%         \[
%         \bY^i_{n,*} = \bY^i_n + \sqrt{2\tau} \xi^{i}_{v,n}
%         \]
%         where $\xi^{i}_{v,n} \sim \mathcal{N}(0, \mathbb{I})$.
%     \EndFor

%     \State Interpolate the updated positions $\bX^i_{n,*}$ and $\bY^i_{n,*}$ onto Fourier modes up to truncation $K$ to obtain the Fourier representations $\tilde{\rho}_u^*$ and $\tilde{\rho}_v^*$.

%     \For{$i=1$ to $N_{x,n}$}
%         \State Compute the branching rate $\tilde{r}_u^i = \frac{f_u(Z_u \rho_u(\bX^i_{n,*}), Z_v \rho_v(\bX^i_{n,*}))}{\rho_u(\bX^i_{n,*})}$.
%         \If{$\tilde{r}_u^i > 0$}
%             \State Duplicate $\bX^i_{n,*}$ with probability $1 - \exp(-\tilde{r}_u^i \tau)$.
%         \ElsIf{$\tilde{r}_u^i < 0$}
%             \State Kill $\bX^i_{n,*}$ with probability $1 - \exp(\tilde{r}_u^i \tau)$.
%         \EndIf
%     \EndFor

%     \For{$i=1$ to $N_{y,n}$}
%         \State Compute the branching rate $\tilde{r}_v^i = \frac{f_v(Z_u \rho_u(\bY^i_{n,*}), Z_v \rho_v(\bY^i_{n,*}))}{\rho_v(\bY^i_{n,*})}$.
%         \If{$\tilde{r}_v^i > 0$}
%             \State Duplicate $\bY^i_{n,*}$ with probability $1 - \exp(-\tilde{r}_v^i \tau)$.
%         \ElsIf{$\tilde{r}_v^i < 0$}
%             \State Kill $\bY^i_{n,*}$ with probability $1 - \exp(\tilde{r}_v^i \tau)$.
%         \EndIf
%     \EndFor

%     \State Update the population sizes: $N_{x,n+1}$ and $N_{y,n+1}$.
%     \State Increment the step: $n = n + 1$.
% \EndWhile
% \end{algorithmic}
% \end{algorithm}

% \section{Kolmogorov forward equation and }
% Recall (overdamped) Langevin diffusion is the solution of the following stochastic differential equation,
% where $\bX_t \in \mathbb R^d$ and .  The probability density function $\rho(t)$ of \eqref{equ:overd_langevin} solves the Kolmogorov forward equation on $\mathbb R^d$
% \begin{equation} \label{equ:FPE}
%     \partial_t  \rho =  \nabla \cdot ( \nabla \rho +\rho \bF).
% \end{equation}
% The main idea of the work is to use the branching process to approximate the evolution of the nonlinearity locally. Before combing with the Langevin dynamics, let us consider the pure branching process given by 
% \begin{equation}
%     \partial_t \rho =  V(x) \rho. 
% \end{equation}
% It can be viewed as an infinite system of ordinary differential equations with index $x\in \mathbb R^d$ without any spatial derivative. 
% If $V(t,x) > 0$, then particles located at $x$ undergo duplication, representing local exponential growth. Conversely, if $V(t,x)<0$, particles at 
% $x$ are removed, corresponding to exponential decay. In this way, the pure branching process encodes a pointwise multiplicative effect on the density, without involving any spatial interactions.
% Combining the Langevin dynamics with the branching process yields the following equation for the evolution of the probability density: 
% \begin{equation} \label{equ:FPE_BDL}
% \partial_t \rho = \nabla \cdot ( \nabla \rho + \rho \bF ) + V(x) \rho, 
% \end{equation} 
% where the first term represents diffusion and drift due to the Langevin process, and the second term captures local growth or decay of mass driven by the branching mechanism.

% The solution of \eqref{equ:FPE_BDL} can be approximated by the empirical measure $\mu^{N}_t$ of a collection of interacting particles $\{x_t^i\}_{i=1}^{N_t}$ evolving as followings $\mu^{N}_t = \frac{1}{N} \sum_{i=1}^{N_t} \delta_{x_t^i} $:
% \begin{enumerate}
%     \item[Step 1:] Between birth/death events, each particle $\bX^i_t$ diffuses independently according to \eqref{equ:overd_langevin}.
%     \item[Step 2:] Each particle also has an independent instantaneous birth-death rate $V(\bX_t^i).$
% \end{enumerate}
% Therefore, we can approximate the solution of \eqref{equ:FPE_BDL} by a density estimation.  Unlike the classical mesh-based method, this approach does not require spatial discretization. This is the key reason why this method is promising to overcome the curse of dimensionality in high-dimensional PDEs.


% \begin{proposition}
% Let $\mu^N_t$ be the empirical measure of particles defined above. Assume that $\mu_0^N \rightharpoonup \rho_0$ as $N\to \infty$. Then for all $t\in (0,\infty)$, we have $\mu^N_t \to \rho_t$, where $\rho_t$ solves \eqref{equ:FPE_BDL} with initial condition $\rho_0$
% \end{proposition}
% \begin{proof}
%     First, we notice that the evolution of $\mu^N_t$ is Markovian since the particles $\bX^i_t$ are, and these particles are interchangeable. Following the proof in \cite{rotskoff2019global}, we want to derive the generator and the backward Kolmogorov equation of the measure-valued Markov process. Let's define for some smooth functional $\Psi:\mathcal P  (\mathbb R ^d) \to \mathbb R $, the generator $\mathcal L_N$
%     \[
%     (\mathcal L_N \Psi ) (\mu^N) := \lim_{t\to 0^+} \frac{\mathbb E_0 \Psi ( \mu_t^N) - \Psi(\mu^N)}{t},
%     \]
%     where $\mathbb E_0 $ denotes the expectation of $\Psi(\mu^N_t)$ conditioned on that $\mu_0^N=\mu^N$. To compute the limit, notice that if particle $\bX_t^i$ is killed at time $t$, the changes this indeuced on $\mu^N_t$ is 
%     \[
%     \mu^N_t =  \mu_{t_-}^N - \frac{1}{N} \delta_{\bX^i}.
%     \]
%     Similarly if particle $\bX_t^i$ is duplicated at time $t$, the changes this indeuced on $\mu^N_t$ is 
%     \[
%     \mu^N_t =  \mu_{t_-}^N + \frac{1}{N} \delta_{\bX^i}.
%     \]
%     If we account for the rate at which these events occur, as well as the effect of the diffusion process, we can explicitly compute the generator 
%     \[
%     \begin{aligned}
%     (\mathcal L _N \Psi)(\mu^N) =& \frac{1}{N} \sum_i ^N \int \left( 
%     \bF(\bX^i) \cdot \nabla D_{\mu^N} \Psi(\bX^i)  
%     + \Delta D_{\mu^N} \Psi(\bX^i)
%     \right) \delta_{\bX^i}(\intd x) \\
%     &+ \frac{1}{N} \sum_{i=1}^N \int V(\bX^i) \delta_{\bX^i}(\intd x)\left(\Psi(\mu^N + \frac{1}{N} \text{sign} V(\bX^i)\delta_{\bX^i} ) - \Psi(\mu^N)\right),
%     \end{aligned}
%     \]
%     where the functional derivative $D_\mu \Psi$ is the function from $\mathbb R^d \to \mathbb R$ defined by that for any signed measure $\nu$ with $\int \nu(\intd x)= 0  $, 
%     \[
%     \lim_{\epsilon \to 0 } \frac{\Psi(\mu+\epsilon\nu)-\Psi(\mu)}{\epsilon} = \int D_\mu\Psi(x)\nu(\intd x).
%     \]
%     By the definition of the empirical measure $\mu^N$, the generator $\mathcal L_N$ can be rewritten as 
%     \[
%     \begin{aligned}
%     (\mathcal L _N \Psi)(\mu^N) =& \int \left( 
%     \bF(x) \cdot \nabla D_{\mu^N} \Psi(x)  
%     + \Delta D_{\mu^N} \Psi(x)
%     \right) \mu^N(\intd x)  \\
%     +&  \int V(x) \mu^N(\intd x) \left(\Psi(\mu^N + \frac{1}{N} \text{sign} V(x)\delta_{x} ) - \Psi(\mu^N)\right) ,
%     \end{aligned}
%     \]
% \end{proof}
% \section{Stochastic Branching method}
% In this work, we aim to extend the above framework to fully nonlinear PDEs. A natural strategy is to adopt an iterative scheme that treats the nonlinear equation as a linear one at each iteration, conditioned on the current estimate of the solution. Specifically, at each time $t$, the equation \eqref{equ:KS} can be approximated by the following linearized system: 
% \begin{equation}
% \begin{aligned}
%     \partial_t \rho_u &=  \nabla \cdot \left ( \bF_u(t,x) \rho_u +\nabla \rho_u \right) + V_u(t,x) \rho_u,\\
%     \partial_t \rho_v &=  \nabla \cdot \left ( \bF_v(t,x) \rho_v +\nabla \rho_v \right) + V_v(t,x) \rho_v,
% \end{aligned}
% \end{equation}
% where the drift and reaction terms are given by $\bF_u =  -\frac{\chi(u)}{u}\nabla v, \bF_v = 0, V_u = \frac{f_u(u,v)}{u}$ and $V_v= \frac{f_v(u,v)}{v}$. At each iteration, the fields $u$ and $v$ are reconstructed from the particle densities as  $u = C_u \rho_u$ and $v= C_v\rho_v$, where $C_u = \int u(0,\bx)\intd \bx$ and $C_v = \int v(0,\bx)\intd \bx$ are normalization constants.

\section{Numerical Examples}
In this section, we present numerical examples in two-dimensional space to verify the accuracy of our method.


\textbf{Case 1:} We begin by solving the following problem in a periodic domain $\Omega = [0, 2\pi] \times [0, 2\pi]$:
\begin{equation} \label{equ:case1}
\begin{aligned}
&u_t - \nabla \cdot \left(\nabla u - u \nabla v \right) = 0  , &\quad \bx \in \Omega , t > 0, \\
&v_t - \Delta v = u - v , & \quad \bx \in \Omega, t > 0.
\end{aligned}
\end{equation}
The system is subject to periodic boundary conditions, with initial conditions given by:
$$
u_0 (\bx) =  \sin^2(\bx_1) \cos^2(\bx_2), \quad v_0(\bx) =  \cos(\bx_1) + \cos(\bx_2) + 2.
$$
The temporal evolution of the total mass of $v$ is characterized by
$$
M(t) = \int v(t,\bx) \, \intd \bx.
$$
It can be shown that the rate of change of the mass is given by
$$
\dot{M}(t) = -M(t) + \int u(x,0) \, dx = -M(t) + \pi^2.
$$
We present the numerical results for $M(t)$, comparing our method with the exact solution in Figure~\ref{fig:total_mass}.

\begin{figure}[h]
\centering
\includegraphics[width=0.8\linewidth]{figure/total_mass.pdf}
\caption{Time evolution of the total mass of $v$ for different numbers of particles. The left panel shows the total mass over time, while the right panel displays the corresponding error in the total mass.}
\label{fig:total_mass}
\end{figure}

Additionally, we compare our particle-based method with finite difference method with a grid of $100 \times 100$ points. The result for $u(t)$ and $v(t)$ at $t = 0.0, 0.1, 0.2, 0.3, 0.4$ are shown in Figures~\ref{fig:fd_comparison_u} and~\ref{fig:fd_comparison_v}, respectively. Figure~\ref{fig:convergence_error} illustrates the reduction in relative $L^2$ error as the number of samples increases, showing a clear $1/\sqrt{N}$ convergence rate.

\begin{figure}[h]
\centering
\includegraphics[width=\linewidth]{figure/case1_u.pdf}
\caption{Comparison of $u$ between the finite difference method with a $100 \times 100$ grid and our particle-based method at $t = 0.0, 0.1, 0.2, 0.3, 0.4$.}
\label{fig:fd_comparison_u}
\end{figure}

\begin{figure}[h]
\centering
\includegraphics[width=\linewidth]{figure/case1_v.pdf}
\caption{Comparison of $v$ between the finite difference method with a $100 \times 100$ grid and our particle-based method at $t = 0.0, 0.1, 0.2, 0.3, 0.4$.}
\label{fig:fd_comparison_v}
\end{figure}

\begin{figure}[h]
\centering
\includegraphics[width=0.8\linewidth]{figure/case1_err.pdf}
\caption{Relative $L^2$ error of $u$ and $v$ with respect to the number of samples $N_0$, compared with $N_0 = 320,000$.}
\label{fig:convergence_error}
\end{figure}

\textbf{Case 2:} For the second case, we solve the system \eqref{equ:case1} with the initial conditions:

$$
u_0(\bx) = 840 \exp\left(-84 (\bx_1^2 + \bx_2^2)\right), \quad v_0(\bx) = 420 \exp\left(-42(\bx_1^2 + \bx_2^2)\right).
$$

The result for $u$ is shown in Figure~\ref{fig:case2_u}, where we observe that the solution undergoes a blow-up, a phenomenon that is generally difficult to capture using traditional function spaces. Furthermore, the $L^2$ norm of the numerical solution for different values of $N_0$ is shown in Figure~\ref{fig:case2_l2}, demonstrating that varying $N_0$ results in similar growth patterns, in contrast to the discontinuous Galerkin method \cite{li2017local}, where the solution is highly sensitive to the mesh size.

\begin{figure}
\centering
\includegraphics[width=0.8\linewidth]{figure/case2_u.pdf}
\caption{Numerical approximations of $u$ at $t = 0, 5 \times 10^{-5}, 10^{-4}, 1.5 \times 10^{-4}$.}
\label{fig:case2_u}
\end{figure}

\begin{figure}
\centering
\includegraphics[width=0.5\linewidth]{figure/case2_L2.pdf}
\caption{Time evolution of the $L^2$-norm of the numerical approximation for different values of $N_0$.}
\label{fig:case2_l2}
\end{figure}

\textbf{Case 3:} Lastly, we examine a fully nonlinear case with logistic growth for the cells, where the interaction between the particles is not known a priori. The system is given by:
\begin{equation}
\label{equ:nonliear}
\begin{aligned}
&u_t - \nabla \cdot \left(\nabla u - \frac{4u}{1 + u^2} \nabla v \right) = u(1 - u), & \quad \bx \in \Omega, t > 0, \\
&v_t - \Delta v = u - v, & \quad \bx \in \Omega, t > 0.
\end{aligned}
\end{equation}
The theoretical results of this model can be found in \cite{osaki2002exponential, wang2007classical, wrzosek2004global}. The initial conditions are specified as:
$$
    \begin{aligned}
        u_0(\bx) =&  \exp\left( (\bx_1 - 0.5\pi)^2 + (\bx_2 - 0.5\pi)^2 \right) + \exp\left( (\bx_1 - \pi)^2 + (\bx_2 - \pi)^2 \right) \\
        &+ \exp\left( (\bx_1 - 1.5\pi)^2 + (\bx_2 - 1.5\pi)^2 \right), \\
        v_0(\bx) =&  \exp\left( (\bx_1 - 0.4\pi)^2 + (\bx_2 - 0.4\pi)^2 \right) + \exp\left( (\bx_1 - 0.8\pi)^2 + (\bx_2 - 0.8\pi)^2 \right) \\
        &+ \exp\left( (\bx_1 - 1.2\pi)^2 + (\bx_2 - 1.2\pi)^2 \right) + \exp\left( (\bx_1 - 1.6\pi)^2 + (\bx_2 - 1.6\pi)^2 \right).
    \end{aligned}
$$
Since both the birth and death of cells, as well as the chemical concentrations $f_u$ and $f_v$, are non-zero, the particle method proposed in \cite{wang2025novel}, which only accounts for advection and diffusion, has limited capability in addressing these effects.  In contrast, our method is sufficiently general to handle these additional complexities.
Figures~\ref{fig:case3_u} and~\ref{fig:case3_v} display the time evolution of the solutions $u$ and $v$, respectively, obtained using both the particle-based and finite difference methods.

\begin{figure}
\centering
\includegraphics[width=0.8\linewidth]{figure/case3_u.pdf}
\caption{Time evolution of the solution $u$ computed using the particle-based and finite difference methods.}
\label{fig:case3_u}
\end{figure}

\begin{figure}
\centering
\includegraphics[width=0.8\linewidth]{figure/case3_v.pdf}
\caption{Time evolution of the solution $v$ computed using the particle-based and finite difference methods.}
\label{fig:case3_v}
\end{figure}

\section{Discussion}
The framework presented here offers a fresh particle-based approach to solving reaction–diffusion--advection equations. Unlike traditional particle or Lagrangian schemes, our method is not restricted to conservation laws and is able to capture aggregation and focusing phenomena in the density variable with high fidelity. This highlights its potential as a useful tool for studying a broad class of nonlinear transport and interaction problems.

From the theoretical perspective, the current work provides heuristic intuition primarily in the linear case. Developing a rigorous convergence theory in more general nonlinear settings remains a challenging but exciting direction for future research. On the numerical side, while the interpolation of measures poses difficulties in high-dimensional problems, this limitation also points to promising avenues for methodological innovation. In particular, Fourier interpolation combined with low-rank approximations~\cite{tang2025solving} could provide a powerful mechanism for scaling the algorithm efficiently to higher dimensions.
% \section{Example Section}
% \label{sec1}
% %% Labels are used to cross-reference an item using \ref command.

% Section text. See Subsection \ref{subsec1}.

% %% Use \subsection commands to start a subsection.
% \subsection{Example Subsection}
% \label{subsec1}

% Subsection text.

% %% Use \subsubsection, \paragraph, \subparagraph commands to 
% %% start 3rd, 4th and 5th level sections.
% %% Refer following link for more details.
% %% https://en.wikibooks.org/wiki/LaTeX/Document_Structure#Sectioning_commands

% \subsubsection{Mathematics}
% %% Inline mathematics is tagged between $ symbols.
% This is an example for the symbol $\alpha$ tagged as inline mathematics.

% %% Displayed equations can be tagged using various environments. 
% %% Single line equations can be tagged using the equation environment.
% \begin{equation}
% f(x) = (x+a)(x+b)
% \end{equation}

% %% Unnumbered equations are tagged using starred versions of the environment.
% %% amsmath package needs to be loaded for the starred version of equation environment.
% \begin{equation*}
% f(x) = (x+a)(x+b)
% \end{equation*}

% %% align or eqnarray environments can be used for multi line equations.
% %% & is used to mark alignment points in equations.
% %% \\ is used to end a row in a multiline equation.
% \begin{align}
%  f(x) &= (x+a)(x+b) \\
%       &= x^2 + (a+b)x + ab
% \end{align}

% \begin{eqnarray}
%  f(x) &=& (x+a)(x+b) \nonumber\\ %% If equation numbering is not needed for a row use \nonumber.
%       &=& x^2 + (a+b)x + ab
% \end{eqnarray}

% %% Unnumbered versions of align and eqnarray
% \begin{align*}
%  f(x) &= (x+a)(x+b) \\
%       &= x^2 + (a+b)x + ab
% \end{align*}

% \begin{eqnarray*}
%  f(x)&=& (x+a)(x+b) \\
%      &=& x^2 + (a+b)x + ab
% \end{eqnarray*}

% %% Refer following link for more details.
% %% https://en.wikibooks.org/wiki/LaTeX/Mathematics
% %% https://en.wikibooks.org/wiki/LaTeX/Advanced_Mathematics

% %% Use a table environment to create tables.
% %% Refer following link for more details.
% %% https://en.wikibooks.org/wiki/LaTeX/Tables
% \begin{table}[t]%% placement specifier
% %% Use tabular environment to tag the tabular data.
% %% https://en.wikibooks.org/wiki/LaTeX/Tables#The_tabular_environment
% \centering%% For centre alignment of tabular.
% \begin{tabular}{l c r}%% Table column specifiers
% %% Tabular cells are separated by &
%   1 & 2 & 3 \\ %% A tabular row ends with \\
%   4 & 5 & 6 \\
%   7 & 8 & 9 \\
% \end{tabular}
% %% Use \caption command for table caption and label.
% \caption{Table Caption}\label{fig1}
% \end{table}


% %% Use figure environment to create figures
% %% Refer following link for more details.
% %% https://en.wikibooks.org/wiki/LaTeX/Floats,_Figures_and_Captions
% \begin{figure}[t]%% placement specifier
% %% Use \includegraphics command to insert graphic files. Place graphics files in 
% %% working directory.
% \centering%% For centre alignment of image.
% \includegraphics{example-image-a}
% %% Use \caption command for figure caption and label.
% \caption{Figure Caption}\label{fig1}
% %% https://en.wikibooks.org/wiki/LaTeX/Importing_Graphics#Importing_external_graphics
% \end{figure}


%% The Appendices part is started with the command \appendix;
%% appendix sections are then done as normal sections
% \appendix
% \section{Example Appendix Section}
% \label{app1}

% Appendix text.

% %% For citations use: 
% %%       \cite{<label>} ==> [1]

% %%
% Example citation, See \cite{lamport94}.

%% If you have bib database file and want bibtex to generate the
%% bibitems, please use
%%
\bibliographystyle{elsarticle-num} 
\bibliography{main.bib}

%% else use the following coding to input the bibitems directly in the
%% TeX file.

%% Refer following link for more details about bibliography and citations.
%% https://en.wikibooks.org/wiki/LaTeX/Bibliography_Management

% \begin{thebibliography}{00}

%% For numbered reference style
%% \bibitem{label}
%% Text of bibliographic item

% \bibitem{lamport94}
%   Leslie Lamport,
%   \textit{\LaTeX: a document preparation system},
%   Addison Wesley, Massachusetts,
%   2nd edition,
%   1994.

% \end{thebibliography}
\end{document}

\endinput
%%
%% End of file `elsarticle-template-num.tex'.
